\section{Proofs for core noise and recourse}
\label{app:recourse}

This appendix proves every result of Section~\ref{sec:recourse}. The
first six subsections build the common machinery: external tools,
finite-law transfer, the exact fallback, certified cells, the all-scale
count, and the projected-growth tail. The remaining subsections prove
the theorems in the order of Section~\ref{sec:recourse}.

\subsection{Conventions and external tools}
\label{app:recourse-tools}

Norms are Euclidean unless marked otherwise; $\norm\cdot_\infty$ is
the maximum norm. $\operatorname{Leb}$ is Lebesgue measure,
$\operatorname{Leb}^*$ its outer measure, $B(c,r)$ the open Euclidean
ball, and $v_k\le2^k$ the volume of the unit ball in $\R^k$. For a
rational number, vector, or matrix $r$, $\langle r\rangle$ is its bit
length (Definition~\ref{def:models-rational}). For a
core-noise instance (Definition~\ref{def:recourse-instance}) we write
$x=(v,z)$, $\kappa_+=\kappa+\sigma>0$, and
\[
 X_C=\{x\in X:\ v\in C\}
\]
for a set $C\subseteq\R^k$. Every instance has a description of $X$ by
at most $L$ explicit rational inequalities: the $2n$ box inequalities in
type~(B), and the rows of $Ex\le e$ together with the supplied box in
type~(P). A bound is \emph{base-computable} if it is computed from the
instance before sampling, in time polynomial in $L$. The letter $b$
denotes the largest bit length of a sampled coefficient. A point of
the grid $G_M$ is $\sigma(2j-M+1)/(M-1)$, so its numerator and
denominator are those of $\sigma$ multiplied by integers of absolute
value less than $M$; since $M$ is a power of two, $b\le L+2\log_2M$. We
use the following tools.

\paragraph{Convex values.}
Lemma~\ref{lem:convex-value} supplies the convex value interface. We
use it in this form. Let $Q$ be a polyhedron given by explicit rational
inequalities and contained in a supplied rational box, and let $\phi$ be
an explicitly encoded rational polynomial of degree at most $D$ that is
convex on $Q$. Given a rational $\eta>0$, one decides whether $Q$ is
empty and, if not, computes a rational $x\in Q$ and a rational $\ell$
with
\begin{equation}
 \ell\le\min_Q\phi\le\phi(x)\le\ell+\eta,
 \label{eq:recourse-cv}
\end{equation}
in bit work $P_D(L_Q+\langle\eta\rangle)$, where $L_Q$ is the bit length
of $Q$, its box, and $\phi$. The tolerance enters through its bit
length $\langle\eta\rangle$, not only through $\log(1/\eta)$: a rational
close to $1$ can have long numerator and denominator. Every tolerance
used below has bit length polynomial in $L$ and the query precision.
Rational linear programs, including
feasibility problems, are solved exactly in polynomial time
\citep[Theorem~(6.4.12) and Section~6.5]{GroetschelLovaszSchrijver1988},
and quadratic programs on nonempty bounded rational polytopes with
rational data and a positive semidefinite objective Hessian have
rational optimal solutions computable in polynomial time
\citep{KozlovTarasovKhachiyan1980}.

\paragraph{Block quantifier elimination.}
We use formulas
\begin{equation}
 \Phi(t)=Q_1x^{(1)}\,Q_2x^{(2)}\ \Psi\bigl(t,x^{(1)},x^{(2)}\bigr),
 \qquad Q_1,Q_2\in\{\exists,\forall\},
 \label{eq:recourse-two-block}
\end{equation}
with one free variable $t$, blocks of $n_1$ and $n_2$ real variables,
and a Boolean combination $\Psi$ of at most $s$ atoms $g\vartriangleright0$,
$\vartriangleright\in\{<,=,>\}$, whose polynomials have degree at most
$d\ge2$. Put
\begin{equation}
 H(s,d,n_1,n_2)=(sd)^{c_{\rm QE}(n_1+1)(n_2+1)}.
 \label{eq:recourse-H}
\end{equation}

\begin{theorem}[Block quantifier elimination]
\label{thm:recourse-qe}
There are an algorithm and an explicit absolute constant $c_{\rm QE}$
with the following properties, where $H$ is given by
\eqref{eq:recourse-H}.
\begin{enumerate}[label=(\alph*)]
\item For arbitrary real coefficients, $\Phi(t)$ is equivalent to a
quantifier-free formula in $t$ that is a Boolean combination of at most
$H$ atoms $h(t)\vartriangleright0$ with $\deg h\le H$.
\item If all coefficients are integers of bit length at most
$\tau\ge2$, the algorithm computes such a formula, with integer
coefficients of bit length at most $(\tau+1)H$, in at most
$\tau^{c_{\rm QE}}H$ bit operations.
\end{enumerate}
\end{theorem}

This is the case of one free variable and two quantifier blocks of
Renegar's theorem \citep[Theorem~1.1]{Renegar1992QE}, whose bounds have
the form $(sd)^{2^{O(\omega)}\ell\prod_jn_j}$ for $\omega$ blocks and
$\ell$ free variables, with bit cost of order
$\tau\log\tau\log\log\tau$ times the same factor and output coefficient
length linear in $\tau$ times that factor; see also
\citet[Chapter~14]{BasuPollackRoy2006}. Part~(a) uses the statement of
that theorem in the real-number model, whose coefficients are arbitrary
reals; it is applied below with noise values and thresholds as
coefficients. Part~(b) uses the separate statement for integer
coefficients in the bit model, and it is applied only after denominators
are cleared. Two features matter below.
The format factor $H$ does not depend on the sizes of the coefficients,
and the coefficient length $\tau$ enters the bit cost and output length
only polynomially, with an absolute exponent. A general
cylindrical-decomposition bound, doubly exponential in the number of
variables, would not suffice.

\begin{theorem}[Univariate real algebraic numbers]
\label{thm:recourse-univariate}
There is an absolute constant $c_U$ with the following property. Given
nonzero integer polynomials $P,P_1,\ldots,P_N$ in one variable, of
degree at most $\delta$ and coefficient bit length at most $\tau$, and
an integer $q\ge0$, one computes the squarefree part of $P$, disjoint
rational isolating intervals for all real roots of $P$, the sign of
every $P_j$ at every such root, and refinements of all isolating
intervals to width at most $2^{-q}$, in $(N+\delta+\tau+q)^{c_U}$ bit
operations \citep[Chapters~8 and~10]{BasuPollackRoy2006}.
\end{theorem}

\begin{theorem}[Lattice reduction]
\label{thm:recourse-lll}
Given a basis of a lattice $\mathcal L\subseteq\Z^{s+1}$ of rank $s$
with entries of bit length at most $\tau$, the LLL algorithm with
parameter $3/4$ computes, in time polynomial in $s+\tau$, a basis whose
first vector $\mathbf b_1$ satisfies
$\norm{\mathbf b_1}\le2^{(s-1)/2}\lambda_1(\mathcal L)$, where
$\lambda_1$ is the length of a shortest nonzero lattice vector
\citep{LenstraLenstraLovasz1982}.
\end{theorem}

\paragraph{Classical analytic facts.}
We use: (F1) a finite convex function on $\R^k$ is differentiable almost
everywhere, and it is differentiable at every point where its
subdifferential is a singleton; a finite concave function has a
nonempty superdifferential at every interior point of its domain; and
$y\in\partial\phi(c)$ if and only if $\phi(c)+\phi^*(y)=y^{\mathsf T}c$
\citep[Theorems~23.4, 23.5, 25.1, 25.5]{Rockafellar1970}. (F2) If
$S\subseteq\R^k$ and $\varphi:S\to\R^k$ has Lipschitz constant $\ell$,
then $\operatorname{Leb}^*(\varphi(S))\le\ell^k\operatorname{Leb}^*(S)$,
since $k$-dimensional Hausdorff measure equals Lebesgue outer measure on
$\R^k$ and does not increase by more than the factor $\ell^k$ under an
$\ell$-Lipschitz map; the covering proof applies to maps defined only
on $S$ \citep[Section~2.4]{EvansGariepy2015}. (F3) If $\phi$ is
convex and differentiable on $\R^k$ with $\kappa$-Lipschitz gradient,
then $\phi(y+u)\le\phi(y)+\nabla\phi(y)^{\mathsf T}u+\frac\kappa2\norm u^2$,
and, if $\inf\phi>-\infty$,
$\phi(y)-\inf\phi\ge\norm{\nabla\phi(y)}^2/(2\kappa)$; the second
inequality follows from the first with $u=-\nabla\phi(y)/\kappa$.

\subsection{Finite laws and one-dimensional sections}
\label{app:recourse-sections}

\begin{lemma}[Section counts]
\label{lem:recourse-sections}
Let $\Phi(t)$ be a formula \eqref{eq:recourse-two-block} with arbitrary
real coefficients, and let $H$ be as in \eqref{eq:recourse-H}. Then
$\{t\in\R:\Phi(t)\}$ and its complement are each a union of at most
$2H^2+1$ intervals, some of which may be single points.
\end{lemma}

\begin{proof}
By Theorem~\ref{thm:recourse-qe}(a), $\Phi$ is a Boolean combination of
at most $H$ atoms $h(t)\vartriangleright0$ with $\deg h\le H$. An
identically zero polynomial has constant sign. The nonzero polynomials
have at most $H^2$ real roots in total. Between consecutive roots, and
at each root, every sign is constant, so the truth value of $\Phi$ is
constant on each of the at most $H^2+1$ open intervals and $H^2$ points
determined by the roots. Both sets are unions of some of these pieces.
\end{proof}

\begin{lemma}[Replacing continuous marginals by grid marginals]
\label{lem:recourse-replacement}
Let $E\subseteq\R^k$ be a Borel set such that, for every coordinate
$j$ and every fixed value of the other $k-1$ coordinates, the section
$\{t:(\ldots,t,\ldots)\in E\}$ is a union of at most $C$ intervals or
points. Let $\Pr_{\rm cont}$ be the uniform law on $[-\sigma,\sigma]^k$
and $\Pr_{\rm grid}$ the uniform law on $G_M^k$. Then
$\abs{\Pr_{\rm grid}(E)-\Pr_{\rm cont}(E)}\le2kC/M$.
\end{lemma}

\begin{proof}
Let $F_c(t)=(t+\sigma)/(2\sigma)$ on $[-\sigma,\sigma]$ and let $F_g$
be the distribution function of the uniform law on $G_M$. For the grid
points $t_j=-\sigma+2\sigma j/(M-1)$ and
$t_j\le t<t_{j+1}$ we have $F_g(t)=(j+1)/M$ and
$j/(M-1)\le F_c(t)<(j+1)/(M-1)$. Hence
$-\frac{j+1}{M(M-1)}\le F_g(t)-F_c(t)\le\frac{M-1-j}{M(M-1)}$, so
$\abs{F_g-F_c}\le1/M$ everywhere; the same holds for left limits, since
$F_c$ is continuous. The probability of an interval or point, with any
convention at its endpoints, is a difference of two such values, so
the two laws differ by at most $2/M$ on it and by at most $2C/M$ on a
union of $C$ of them.

For $0\le j\le k$ let $\nu_j$ be the product law whose first $j$
coordinates are uniform on $G_M$ and whose remaining coordinates are
uniform on $[-\sigma,\sigma]$. Then $\nu_0=\Pr_{\rm cont}$ and
$\nu_k=\Pr_{\rm grid}$. The laws $\nu_{j-1}$ and $\nu_j$ differ only in
coordinate $j$. By Fubini's theorem, $\nu_j(E)-\nu_{j-1}(E)$ is the
average, over the other coordinates distributed according to their
common marginals, of the difference of the two one-dimensional laws on
the section of $E$. Each such difference is at most $2C/M$ in absolute
value, whatever the values of the other coordinates. Summing over $j$
proves the claim.
\end{proof}

Every event to which we apply Lemma~\ref{lem:recourse-replacement} is
defined by a formula \eqref{eq:recourse-two-block} in which one noise
coordinate is the free variable and all other noise coordinates and
thresholds are real coefficients. It is semialgebraic, hence Borel, and
Lemma~\ref{lem:recourse-sections} gives the section bound. The bound
covers ties and exact algebraic relations among the noise values; no
almost-sure statement replaces it.

\subsection{The selected point and the exact fallback}
\label{app:recourse-fallback}

\begin{lemma}[The selected point]
\label{lem:recourse-selected-point}
Let $(X,f,\kappa,\sigma)$ be a core-noise instance and $\gamma\in\R^k$.
\begin{enumerate}[label=(\alph*)]
\item The set $\argmin_Xf_\gamma$ is nonempty and compact, and the
lexicographically least core $a$ among its points exists.
\item The fiber $S(\gamma)$ is nonempty, compact, and convex. Its point
$p$ of least norm is unique, and $x^\sharp=(a,p)$ is a global minimizer.
\item Let
$\operatorname{LexGT}(v',v)$ be the formula
$(v'_1>v_1)\vee\bigl(v'_1=v_1\wedge\bigl((v'_2>v_2)\vee(v'_2=v_2\wedge
(\cdots\vee(v'_{k-1}=v_{k-1}\wedge v'_k>v_k)\cdots))\bigr)\bigr)$, with
$2k-1$ atoms, and put
\[
 \operatorname{Sel}(x',x)=\operatorname{LexGT}(v',v)\vee
 \bigl(v'=v\wedge\norm{z'}^2\ge\norm z^2\bigr).
\]
For each coordinate index $i$ of $\R^n$, the formula
\begin{equation}
 \exists x\ \forall x':\ x\in X\wedge x_i=t\wedge
 \bigl[x'\notin X\vee f_\gamma(x')>f_\gamma(x)\vee
 \bigl(f_\gamma(x')=f_\gamma(x)\wedge\operatorname{Sel}(x',x)\bigr)\bigr]
 \label{eq:recourse-singleton}
\end{equation}
holds exactly for $t=x^\sharp_i$.
\end{enumerate}
\end{lemma}

\begin{proof}
(a) $X$ is compact and $f_\gamma$ is continuous. Minimizing $v_1$ over
the compact optimal set, then $v_2$ over the compact set of minimizers,
and so on, produces the lexicographically least core.

(b) In type~(B), $S(\gamma)=\argmin_{z\in[0,1]^m}f_\gamma(a,z)$, because
$a$ is an optimal core; this set is convex because $f(a,\cdot)$ is
convex on $[0,1]^m$. In type~(P), let $X_a=\{z:(a,z)\in X\}$. The
function $z\mapsto f(a,z)+\frac\alpha2\norm a^2$ is the restriction of
the convex function $f+\frac\alpha2\norm v^2$ to an affine slice of $X$,
so $f_\gamma(a,\cdot)$ is convex on $X_a$ and
$S(\gamma)=\argmin_{X_a}f_\gamma(a,\cdot)$ is convex. Compactness and
nonemptiness are clear, and the least-norm point of a nonempty compact
convex set is unique.

(c) Let $x$ satisfy the matrix of \eqref{eq:recourse-singleton}. Every
feasible $x'$ has $f_\gamma(x')\ge f_\gamma(x)$, so $x$ is a global
minimizer. If its core $v$ differed from $a$, then $a$ would be
lexicographically smaller, and a global minimizer $x'$ with core $a$
would violate the bracket. Hence $v=a$, $z\in S(\gamma)$, and the
competitors $(a,z')$ with $z'\in S(\gamma)$ force
$\norm z\le\norm{z'}$; thus $z=p$. Conversely $x^\sharp$ satisfies the
matrix: a competitor of equal value is a global minimizer, so its core
is lexicographically at least $a$, and if its core equals $a$ its
residual lies in $S(\gamma)$ and has norm at least $\norm p$.
\end{proof}

\begin{lemma}[Exact fallback with a base-only exponential factor]
\label{lem:recourse-fallback}
Fix $D$. There are an absolute constant $c_F$ and a polynomial $p_D$
such that every core-noise instance has a base-computable integer
$\Xi_0\le2^{p_D(L)}$ with the following property. For every rational
$\gamma\in[-\sigma,\sigma]^k$ whose entries have bit length at most
$b$, and every integer $q\ge0$, one computes a rational $\tilde x\in X$
with $\norm{\tilde x-x^\sharp(\gamma)}_\infty\le2^{-q}$ and bit length
at most $(L+b+q+1)^{c_F}$, using at most $\Xi_0(L+b+q+1)^{c_F}$ bit
operations. Consequently, with a base-computable rational
$G\ge\max\{1,\sup_X\norm{\nabla f_\gamma}\}$ valid for all
$\gamma\in[-\sigma,\sigma]^k$, one obtains within the same bound a
rational $x_q\in X$ and $\ell_q$ satisfying
\eqref{eq:recourse-core-output}.
\end{lemma}

\begin{proof}
Clear the denominators of \eqref{eq:recourse-singleton}. Its polynomials
have degree at most $d=\max\{D,2\}$ and integer coefficients of bit
length $\tau\le c(L+kb)$ for an absolute $c$, because the products of
all denominators of $f$, of the description of $X$, and of $\gamma$ have
total length at most $L+kb$. The formula has two blocks of $n$
variables, one free variable, and $s\le4L+4n+2$ atoms. Apply
Theorem~\ref{thm:recourse-qe}(b) to each of the $n$ formulas. With
$H=H(s,d,n,n)$, the cost is at most $n\tau^{c_{\rm QE}}H$, and each
output formula has at most $H$ atoms of degree at most $H$ with
coefficients of bit length at most $(\tau+1)H$.

Fix $i$ and the output formula $\Phi_i(t)$. By
Lemma~\ref{lem:recourse-selected-point}(c) its solution set is
$\{x^\sharp_i\}$. Some atom polynomial of $\Phi_i$ is nonconstant and
vanishes at $x^\sharp_i$: otherwise every atom would keep its sign on a
neighborhood of $x^\sharp_i$, and $\Phi_i$ would hold on that
neighborhood. Let $P$ be the product of the nonconstant atom
polynomials; its degree is at most $H^2$ and its coefficient length is
at most $H^2((\tau+1)H+\log(H+1))$. Theorem~\ref{thm:recourse-univariate}
isolates the real roots of the squarefree part of $P$ and determines the
signs of all atom polynomials at each root; evaluating the Boolean
combination identifies the unique root satisfying $\Phi_i$. Refining
its isolating interval to width $2^{-q-2}$ and rounding the midpoint to
a dyadic number with $q+3$ fractional bits gives $w_i$ with
$\abs{w_i-x^\sharp_i}\le2^{-q-1}$; since $\abs{x^\sharp_i}$ is bounded
by the supplied box, $w_i$ has $O(q+L)$ bits. These steps cost
$(H^3+\tau H^3+q)^{c_U}$ per coordinate.

In type~(B), clip $w$ to the box: the result $\tilde x$ is feasible and
no farther from $x^\sharp$ in any coordinate. In type~(P), solve the
rational linear feasibility problem $\{x\in X:\norm{x-w}_\infty\le2^{-q-1}\}$;
it contains $x^\sharp$, and any solution $\tilde x$ satisfies
$\norm{\tilde x-x^\sharp}_\infty\le2^{-q}$ and has bit length polynomial
in $L+q$. Since $\log H$ is polynomial in $L$ for fixed $D$, all costs
are at most $\Xi_0(L+b+q+1)^{c_F}$ for a base-computable
$\Xi_0\le2^{p_D(L)}$ and an absolute $c_F$; here $\Xi_0$ absorbs every
power of $H$ and the factor $n$.

For the last statement, take $G=\max\{1,G_f+k\sigma\}$, where $G_f$
bounds $\norm{\nabla f}$ on the supplied box by coefficient sums and
$k\sigma\ge\norm\gamma$. Apply the first part with
precision $q+1+\lceil\log_2(nG)\rceil=q+O(L)$. Then
$\norm{\tilde x-x^\sharp}\le\sqrt n\,2^{-q-1}/(nG)\le2^{-q}/(2G)$; since $X$
is convex, $\abs{f_\gamma(\tilde x)-f_\gamma^*}\le2^{-q-1}$. Put
$x_q=\tilde x$ and $\ell_q=f_\gamma(\tilde x)-2^{-q-1}$. Then
\eqref{eq:recourse-core-output} holds, and
$\norm{v_q-a}\le\norm{\tilde x-x^\sharp}\le2^{-q}$.
\end{proof}

The base factor $\Xi_0$ depends only on the format of the instance:
its dimension, the number of inequalities describing $X$, and $D$. The
sampled coefficient length $b$ and the precision $q$ enter only
polynomially. This lets the grid size be chosen after $\Xi_0$, with
$\log M$ polynomial in $L$; if the fallback cost were
$2^{\poly(L+b)}$, the grid choice would be circular.

\subsection{Certified cells}
\label{app:recourse-cells}

For $j\ge0$ and $h=2^{-j}$, a \emph{cell of level $j$} is a box
$C=\prod_{i=1}^k[l_i,l_i+h]$ with $l_i\in h\Z\cap[0,1-h]$; its
\emph{children} are the $2^k$ cells of level $j+1$ inside it. Put
\[
 e_h=\frac{\kappa_+kh^2}8,\qquad \varepsilon_h=4e_h=\frac{\kappa_+kh^2}2 .
\]
A \emph{cell oracle} receives a draw $\gamma$ and a cell $C$ of side $h$.
It returns $\bot$ only if $X_C=\emptyset$; otherwise it returns a
rational number $\operatorname{LB}(C)$ and a rational point
$\omega_C\in X_C$ with
\begin{equation}
 \operatorname{LB}(C)\le f_\gamma(x)\ \ (x\in X_C),\qquad
 f_\gamma(\omega_C)\le\operatorname{LB}(C)+2e_h .
 \label{eq:recourse-cell-oracle}
\end{equation}

\begin{lemma}[Cell oracles]
\label{lem:recourse-cell-oracle}
Instances of both types have cell oracles. For a cell of level $j$ the
oracle uses at most $2^kP_D(L+b+j)$ bit operations, and its outputs
have bit length at most $P_D(L+b+j)$.
\end{lemma}

\begin{proof}
\emph{Type~(B).} For each corner $w$ of $C$, apply
\eqref{eq:recourse-cv} to the convex polynomial $z\mapsto f_\gamma(w,z)$
on $[0,1]^m$ with accuracy $e_h$ (if $m=0$, evaluate $f_\gamma(w)$); it returns $z_w$ and $\ell_w$ with
$\ell_w\le V_\gamma(w)\le f_\gamma(w,z_w)\le\ell_w+e_h$, where
$V_\gamma(w)=\min_zf_\gamma(w,z)$. Put
$\operatorname{LB}(C)=\min_w\ell_w-e_h$ and $\omega_C=(w^*,z_{w^*})$
for a corner $w^*$ attaining the minimum. The second inequality in
\eqref{eq:recourse-cell-oracle} is immediate. For the first, let
$(v,z)\in X_C$. Choose independent random corner coordinates
$W_i\in\{l_i,l_i+h\}$ with $\Pr\{W_i=l_i+h\}=(v_i-l_i)/h$, so that
$\E W_i=v_i$ and $\operatorname{Var}W_i\le h^2/4$. Replace the core
coordinates one at a time. Conditionally on $W_1,\ldots,W_{i-1}$, the
function $\varphi(u)=f_\gamma(W_1,\ldots,W_{i-1},u,v_{i+1},\ldots,v_k,z)$
satisfies $\varphi''\le\Lambda$ on $[l_i,l_i+h]$, so
$\varphi(u)\le\varphi(v_i)+\varphi'(v_i)(u-v_i)+\frac\Lambda2(u-v_i)^2$
there, and $\E\varphi(W_i)\le\varphi(v_i)+\Lambda h^2/8$. Summing over
$i$, $\E f_\gamma(W,z)\le f_\gamma(v,z)+e_h$, since $\Lambda\le\kappa_+$.
Some corner $w$ therefore has $V_\gamma(w)\le f_\gamma(w,z)\le
f_\gamma(v,z)+e_h$, and
$\operatorname{LB}(C)\le\ell_w-e_h\le f_\gamma(v,z)$.

\emph{Type~(P).} Put
$\phi_C(x)=f_\gamma(x)+\frac{\kappa_+}2\sum_i(v_i-l_i)(v_i-l_i-h)$. Since
$\kappa_+=\alpha+\sigma$, $\phi_C$ equals $f+\frac\alpha2\norm v^2$ plus
the convex function $\frac\sigma2\norm v^2$ plus an affine function, so
it is convex on $X$. On $X_C$ each product $(v_i-l_i)(v_i-l_i-h)$ lies
in $[-h^2/4,0]$, hence $f_\gamma-e_h\le\phi_C\le f_\gamma$ there. Apply
\eqref{eq:recourse-cv} to $\phi_C$ on the polyhedron
$X_C=X\cap\{l\le v\le l+h\}$ with accuracy $e_h$. If $X_C$ is empty,
return $\bot$. Otherwise it returns $\omega_C\in X_C$ and $\ell_C$ with
$\ell_C\le\min_{X_C}\phi_C\le\phi_C(\omega_C)\le\ell_C+e_h$; put
$\operatorname{LB}(C)=\ell_C$. Then $\operatorname{LB}(C)\le
\min_{X_C}f_\gamma$ and $f_\gamma(\omega_C)\le\phi_C(\omega_C)+e_h\le
\operatorname{LB}(C)+2e_h$.

In both types the data of each convex call have bit length
$P_D(L+b+j)$, and the tolerance $e_h=\kappa_+k\,2^{-2j-3}$ is positive,
because $\kappa_+\ge\sigma>0$, and has $\langle e_h\rangle\le2j+P_D(L)$.
\end{proof}

The type~(P) construction is the classical secant underestimator of
$\alpha$BB branch and bound \citep{AndroulakisMaranasFloudas1995}. The
type~(B) construction needs only the diagonal bound $\Lambda$ and convex
residual fibers. We add $\sigma$ to the curvature parameter so that
$e_h>0$ also when $\kappa=0$; this changes $1+\kappa/\sigma$ by at most a
factor two.

\paragraph{Refinement.}
Given a draw, a final level $J$, and a cap $N$, the \emph{refinement}
proceeds as follows. Level $0$ consists of the single cell $[0,1]^k$.
For $j\ge1$, if $2^k$ times the number of cells retained at level $j-1$
exceeds $N$, stop with the outcome \textsc{cap}; otherwise generate all
children of the retained cells. At every level call the cell oracle on
each generated cell, discard the cells that return $\bot$, set the
incumbent $x_{\rm inc}$ to a witness of least value among all witnesses
produced so far and $U=f_\gamma(x_{\rm inc})$, and then, in a second
pass using this final value of $U$, retain exactly the cells with
$\operatorname{LB}(C)\le U$. The list work at a level is linear in the
number of generated cells.

\begin{lemma}[Invariants of the refinement]
\label{lem:recourse-refinement}
After level $j$ has been completed, with $h=2^{-j}$:
\begin{enumerate}[label=(\alph*)]
\item $U-2e_h\le f_\gamma^*\le U=f_\gamma(x_{\rm inc})$ and
$U\le f_\gamma^*+2e_h$;
\item for every global minimizer $x^*=(v^*,z^*)$ of $f_\gamma$, some
retained cell contains $v^*$;
\item some retained cell contains the core of $x_{\rm inc}$;
\item every retained cell $C$ satisfies
$f_\gamma(\omega_C)\le f_\gamma^*+4e_h$ and
$C\subseteq v(\omega_C)+[-h,h]^k$, where $v(\omega_C)$ is the core of
$\omega_C$.
\end{enumerate}
\end{lemma}

\begin{proof}
(b) At level $0$ the single cell contains $v^*$ and its lower bound is
at most $f_\gamma(x^*)=f_\gamma^*\le U$. If a retained cell of level
$j-1$ contains $v^*$, one of its children does; that child is generated,
and \eqref{eq:recourse-cell-oracle} gives
$\operatorname{LB}\le f_\gamma^*\le U$, so it is retained.

(c) The incumbent was produced as a witness $\omega_{C_0}$ of a cell
$C_0$ generated at some level $j_0\le j$, and it has remained the
incumbent since. At level $j_0$, \eqref{eq:recourse-cell-oracle} gives
$\operatorname{LB}(C_0)\le f_\gamma(x_{\rm inc})=U$, so $C_0$ is
retained. Induction over the levels after $j_0$, exactly as in (b) with
$x_{\rm inc}$ in place of $x^*$, proves (c).

(a) The cells retained at level $j$, together with all cells discarded
at levels $0,\ldots,j$, cover $[0,1]^k$; this follows by induction,
since the children of a retained cell cover it. A cell discarded at
level $j'\le j$ either has $X_C=\emptyset$ or has
$\operatorname{LB}(C)>U_{j'}\ge U$, where $U_{j'}$ is the incumbent
value at level $j'$; incumbent values never increase. A cell retained
at level $j$ has $\operatorname{LB}(C)\ge f_\gamma(\omega_C)-2e_h\ge
U-2e_h$, because $\omega_C$ was a candidate for the incumbent. Hence
every feasible point has value at least $U-2e_h$. Finally, the
retained cell containing $v^*$ from (b) has
$f_\gamma(\omega_C)\le\operatorname{LB}(C)+2e_h\le f_\gamma^*+2e_h$,
so $U\le f_\gamma^*+2e_h$.

(d) A retained cell has $f_\gamma(\omega_C)\le\operatorname{LB}(C)+2e_h
\le U+2e_h\le f_\gamma^*+4e_h$ by (a). Since $v(\omega_C)\in C$ and $C$
has side $h$, every point of $C$ differs from $v(\omega_C)$ by at most
$h$ in each coordinate.
\end{proof}

All four invariants hold for every draw and for every admissible answer
of the convex value calls; no growth or probability assumption is used.

\subsection{An all-scale count}
\label{app:recourse-count}

For a draw $\gamma\in\R^k$ and $0<h\le1$ define the projected
near-optimal set, its padded hull, and the count proxy
\begin{align*}
 S_h(\gamma)&=\{v:\ \exists z\ (v,z)\in X,\
  f_\gamma(v,z)\le f_\gamma^*+\varepsilon_h\},\\
 K_h(\gamma)&=\conv S_h(\gamma)+[-h,h]^k,\\
 \mathfrak d(K)&=\max\{\abs{\det(y_1-y_0,\ldots,y_k-y_0)}:\
  y_0,\ldots,y_k\in K\},\\
 A(\gamma)&=\sup_{0<h\le1}\mathfrak d(K_h(\gamma))/h^k,\qquad
 W(\gamma)=\max\{1,A(\gamma)\}.
\end{align*}
The set $S_h(\gamma)$ is nonempty and compact, being the projection of
a compact sublevel set, so $K_h(\gamma)$ is a compact convex body with
nonempty interior. We allow $A(\gamma)=+\infty$.

\begin{lemma}[Simplex volume]
\label{lem:recourse-simplex}
For every compact convex $K\subseteq\R^k$ with nonempty interior,
$\mathfrak d(K)/k!\le\vol(K)\le2^k\mathfrak d(K)$.
\end{lemma}

\begin{proof}
The simplex with vertices attaining $\mathfrak d(K)$ lies in $K$ and has
volume $\mathfrak d(K)/k!$. For the upper bound, fix such vertices
$y_0,\ldots,y_k$; since $\mathfrak d(K)>0$, the edges $y_i-y_0$ form a
basis. Write $y\in K$ as $y=y_0+\sum_i\lambda_i(y_i-y_0)$. Replacing
$y_i$ by $y$ multiplies the determinant by $\lambda_i$, so maximality
gives $\abs{\lambda_i}\le1$. Thus $K$ lies in the parallelepiped
$y_0+\{\sum_i\lambda_i(y_i-y_0):\lambda\in[-1,1]^k\}$ of volume
$2^k\mathfrak d(K)$.
\end{proof}

\begin{lemma}[Cell count]
\label{lem:recourse-cell-count}
For every draw and every run of the refinement, at most $2^kW(\gamma)$
cells are retained at each completed level, and at most $4^kW(\gamma)$
cells are generated at each level.
\end{lemma}

\begin{proof}
Let $h=2^{-j}$. By Lemma~\ref{lem:recourse-refinement}(d), the core of
the witness of each retained cell lies in $S_h(\gamma)$, and the cell
lies in $K_h(\gamma)$. Distinct cells of level $j$ have disjoint
interiors and volume $h^k$. By Lemma~\ref{lem:recourse-simplex} their
number is at most $\vol(K_h)/h^k\le2^k\mathfrak d(K_h)/h^k\le2^kA$. The
next level has at most $2^k$ times as many cells, and level $0$ has one.
\end{proof}

We now bound the tail of $W$. Write $c=-\gamma$, so that
$f_\gamma(x)=f(x)-c^{\mathsf T}v$; the uniform laws on
$[-\sigma,\sigma]^k$ and on $G_M^k$ are symmetric, so $c$ and $\gamma$
have the same law. Define, for $c\in\R^k$,
\[
 \Psi(c)=\max_{x\in X}\bigl(c^{\mathsf T}v-f(x)\bigr)
 =-\min_X f_{-c},\qquad
 \Upsilon(c)=\Psi(c)+\frac{\norm c^2}{2\kappa_+}.
\]

\begin{lemma}[A conjugate pushforward measure]
\label{lem:recourse-conjugate}
\begin{enumerate}[label=(\alph*)]
\item $\Psi$ is finite and convex, and every subgradient of $\Psi$ lies
in $[0,1]^k$. The conjugate $\Upsilon^*$ is finite and differentiable
on $\R^k$, $\nabla\Upsilon^*$ is $\kappa_+$-Lipschitz, and
$y\in\partial\Upsilon(c)$ if and only if $\nabla\Upsilon^*(y)=c$.
\item $\nu(E)=\operatorname{Leb}\{y:\nabla\Upsilon^*(y)\in E\}$ is a Borel
measure, and $\nu(\prod_i[\alpha_i,\beta_i])\le
\prod_i(1+(\beta_i-\alpha_i)/\kappa_+)$.
\item For every $c\in\R^k$ and $0<h\le1$,
$\vol(K_h(-c))\le\nu\bigl(B(c,4\kappa_+kh)\bigr)$.
\end{enumerate}
\end{lemma}

\begin{proof}
(a) $\Psi$ is a maximum of affine functions of $c$ over a compact index
set, hence finite and convex. If $s\in\partial\Psi(c)$ and $u\in\R^k$,
then $s^{\mathsf T}u\le\Psi(c+u)-\Psi(c)\le\sum_i\max\{u_i,0\}$, because
every core lies in $[0,1]^k$. Taking $u=\pm\epsilon e_i$ gives
$0\le s_i\le1$. The function $\Upsilon$ is $(1/\kappa_+)$-strongly
convex. For fixed $y$, the function $c\mapsto y^{\mathsf T}c-\Upsilon(c)$
is strongly concave and continuous, so it has a unique maximizer
$c(y)$ and $\Upsilon^*(y)$ is finite. By (F1),
$\partial\Upsilon^*(y)=\{c:y\in\partial\Upsilon(c)\}$ is the set of
these maximizers, namely $\{c(y)\}$, so $\Upsilon^*$ is differentiable
with $\nabla\Upsilon^*(y)=c(y)$. If $y_j\in\partial\Upsilon(c_j)$ for
$j=1,2$, then $y_j-c_j/\kappa_+\in\partial\Psi(c_j)$, and monotonicity of
$\partial\Psi$ gives
$(y_1-y_2)^{\mathsf T}(c_1-c_2)\ge\norm{c_1-c_2}^2/\kappa_+$; hence
$\norm{c_1-c_2}\le\kappa_+\norm{y_1-y_2}$.

(b) $\nabla\Upsilon^*$ is continuous, so $\nu$ is a Borel measure. If
$\nabla\Upsilon^*(y)=c$, then
$y\in\partial\Upsilon(c)=\partial\Psi(c)+c/\kappa_+\subseteq[0,1]^k+c/\kappa_+$.
For a box $E$ the preimage therefore lies in $[0,1]^k+E/\kappa_+$, a box
with the stated volume.

(c) Let $\gamma=-c$ and $v\in S_h(\gamma)$, with a residual $z$ such that
$(v,z)\in X$ and $f(v,z)-c^{\mathsf T}v\le-\Psi(c)+\varepsilon_h$. Put
$y_0=v+c/\kappa_+$ and $r(y)=\Upsilon(c)+\Upsilon^*(y)-c^{\mathsf T}y\ge0$.
Since $\Psi(c')\ge c'^{\mathsf T}v-f(v,z)$ for all $c'$,
\[
 \Upsilon^*(y_0)\le\sup_{c'}\Bigl(y_0^{\mathsf T}c'-c'^{\mathsf T}v+f(v,z)
 -\frac{\norm{c'}^2}{2\kappa_+}\Bigr)=f(v,z)+\frac{\norm c^2}{2\kappa_+},
\]
and therefore $r(y_0)\le\Psi(c)+f(v,z)-c^{\mathsf T}v\le\varepsilon_h$.
The function $r$ is convex, has $\kappa_+$-Lipschitz gradient
$\nabla\Upsilon^*-c$, and has infimum zero. By (F3),
$\norm{\nabla r(y_0)}\le\sqrt{2\kappa_+\varepsilon_h}=\kappa_+\sqrt kh$,
and for $u\in[-h,h]^k$, where $\norm u\le\sqrt kh$,
\[
 r(y_0+u)\le\varepsilon_h+\kappa_+\sqrt kh\norm u+\frac{\kappa_+}2\norm u^2
 \le2\kappa_+kh^2.
\]
The sublevel set $\{r\le2\kappa_+kh^2\}$ is convex and contains
$S_h(\gamma)+c/\kappa_+ +[-h,h]^k$, hence also
$K_h(\gamma)+c/\kappa_+$. On it, (F3) gives
$\norm{\nabla\Upsilon^*(y)-c}=\norm{\nabla r(y)}\le2\kappa_+\sqrt kh<4\kappa_+kh$.
Thus $K_h(\gamma)+c/\kappa_+\subseteq\{y:\nabla\Upsilon^*(y)\in
B(c,4\kappa_+kh)\}$, and translation invariance of volume proves (c).
\end{proof}

\begin{lemma}[Disjoint balls]
\label{lem:recourse-vitali}
Let $\mathcal F$ be a family of open balls $B(c,r_c)$, one for each
center $c$ in a set $E\subseteq\R^k$, with $0<r_c\le R_*$, all contained
in a bounded set. There is a countable disjoint subfamily
$\{B(c_j,r_j)\}$ with $E\subseteq\bigcup_jB(c_j,5r_j)$.
\end{lemma}

\begin{proof}
Put $\mathcal F_0=\mathcal F$. While $\mathcal F_i\ne\emptyset$, let
$\rho_i$ be the supremum of the radii in $\mathcal F_i$, choose
$B(c_i,r_i)\in\mathcal F_i$ with $r_i>\rho_i/2$, and let
$\mathcal F_{i+1}$ consist of the balls of $\mathcal F_i$ disjoint from
it. The chosen balls are disjoint. A ball $B(c,r_c)$ removed at step $i$
meets $B(c_i,r_i)$ and has $r_c\le\rho_i<2r_i$, so
$\norm{c-c_i}<r_i+r_c<3r_i$. If some ball $B(c,r_c)$ were never chosen
or removed, then $\rho_i\ge r_c$ for all $i$, and infinitely many
disjoint chosen balls of radius greater than $r_c/2$ would lie in a
bounded set, which is impossible.
\end{proof}

\begin{lemma}[Weak first-moment tail]
\label{lem:recourse-weak-tail}
For every $s>0$,
\[
 \Pr_{\rm cont}\{A>s\}\le\frac{a_k}s,\qquad
 a_k=(180k^3)^k\Bigl(1+\frac{\kappa_+}\sigma\Bigr)^k .
\]
\end{lemma}

\begin{proof}
Let $R_*=4\kappa_+k$, $Q=[-\sigma,\sigma]^k$, and
$E_s=\{c\in Q:A(-c)>s\}$. For $c\in E_s$ choose $h\in(0,1]$ with
$\mathfrak d(K_h(-c))>sh^k$. By Lemmas~\ref{lem:recourse-simplex}
and~\ref{lem:recourse-conjugate}(c), the ball $B(c,r_c)$ with
$r_c=R_*h\le R_*$ satisfies
\[
 \nu(B(c,r_c))\ge\vol(K_h(-c))\ge\frac{\mathfrak d(K_h(-c))}{k!}
 >\frac{s\,r_c^k}{k!R_*^k}.
\]
All these balls lie in $Q_*=[-\sigma-R_*,\sigma+R_*]^k$.
Lemma~\ref{lem:recourse-vitali} selects disjoint balls $B(c_j,r_j)$
whose fivefold enlargements cover $E_s$. Therefore
\[
 \operatorname{Leb}^*(E_s)\le\sum_jv_k5^kr_j^k
 <\frac{5^kv_kk!R_*^k}{s}\sum_j\nu(B(c_j,r_j))
 \le\frac{5^kv_kk!R_*^k}{s}\,\nu(Q_*).
\]
By Lemma~\ref{lem:recourse-conjugate}(b),
$\nu(Q_*)\le(1+2\sigma/\kappa_++8k)^k$. Dividing by $(2\sigma)^k$ and
using $v_k\le2^k$, $k!\le k^k$, and
$\kappa_+/\sigma+2+8k\kappa_+/\sigma\le9k(1+\kappa_+/\sigma)$ gives
\[
 \Pr_{\rm cont}\{A>s\}\le\frac{5^k2^kk^k(2k)^k}{s}
 \Bigl(\frac{\kappa_+}\sigma+2+\frac{8k\kappa_+}\sigma\Bigr)^k
 \le\frac{(180k^3)^k(1+\kappa_+/\sigma)^k}{s}.\qedhere
\]
\end{proof}

The bound has the power $s^{-1}$ in every core dimension. It is the
weak-type estimate for the centered maximal function of the locally
finite measure $\nu$; the conjugate construction converts near-optimal
sets at every scale into balls of comparable $\nu$-mass.

\begin{lemma}[The count event is a two-block formula]
\label{lem:recourse-count-formula}
There is a base-computable $C_W\le2^{P_D(L)}$ such that, for every
power of two $M$ and every $s\ge1$,
$\Pr_{\rm grid}\{W>s\}\le a_k/s+C_W/M$.
\end{lemma}

\begin{proof}
For $s\ge1$, $\{W>s\}=\{A>s\}$, and $A(\gamma)>s$ holds if and only if
there are $h\in(0,1]$ and $y_0,\ldots,y_k\in K_h(\gamma)$ with
$\abs{\det(y_1-y_0,\ldots,y_k-y_0)}>sh^k$. By Carath\'eodory's theorem,
$y_i\in K_h(\gamma)$ if and only if
$y_i=\sum_{l=0}^k\lambda_{il}v_{il}+u_i$ with $\lambda_{il}\ge0$,
$\sum_l\lambda_{il}=1$, $v_{il}\in S_h(\gamma)$, and
$\norm{u_i}_\infty\le h$. Since minima over $X$ are attained,
$v_{il}\in S_h(\gamma)$ if and only if some $z_{il}$ has
$(v_{il},z_{il})\in X$ and
$f_\gamma(v_{il},z_{il})\le f_\gamma(x')+\kappa_+kh^2/2$ for every
$x'\in X$. One universal block $x'$ serves all $(k+1)^2$ comparisons.

The determinant is not expanded into $k!$ monomials. Instead introduce
existential matrices $Q$ and $R$ with
$(y_1-y_0,\ldots,y_k-y_0)=QR$, $Q^{\mathsf T}Q=I$, $R$ upper triangular,
$R_{ii}>0$, and product chains $\pi_1=R_{11}$,
$\pi_i=\pi_{i-1}R_{ii}$, $\epsilon_1=h$, $\epsilon_i=\epsilon_{i-1}h$,
and require $\pi_k>s\,\epsilon_k$. A nonsingular real matrix has such a
factorization, and then $\pi_k$ equals the absolute value of its
determinant; conversely these equations force $\abs{\det}=\pi_k$. The
strict inequality forces nonsingularity.

The resulting formula has the form \eqref{eq:recourse-two-block} with
one existential block of $O(k^2(k+m))$ variables, one universal block
of $n$ variables, $O(k^2(k+L))$ atoms, and degree at most $\max\{D,2\}$;
the free variable is one noise coordinate, and the other noise
coordinates and $s$ are real coefficients. By
Lemma~\ref{lem:recourse-sections}, every section of $\{A>s\}$ is a
union of at most $C=2H^2+1$ intervals or points, with $H$ from
\eqref{eq:recourse-H}, uniformly in $s$ and in the fixed coordinates.
Lemmas~\ref{lem:recourse-replacement} and~\ref{lem:recourse-weak-tail}
give the claim with $C_W=2kC$. Since $\log H$ is polynomial in $L$ for
fixed $D$, so is $\log C_W$.
\end{proof}

The formula also shows that $\{A>s\}$ is semialgebraic, so $A$ and $W$
are measurable, including the event $\{W=\infty\}$, whose probability
is at most $C_W/M$ under the grid law.

\subsection{Projected growth}
\label{app:recourse-growth}

For a draw $\gamma$ let $g(\gamma)$ be the supremum of the numbers
$t\ge0$ for which some $x^*=(v^*,z^*)\in X$ satisfies
\begin{equation}
 f_\gamma(x)-f_\gamma(x^*)\ge t\norm{v-v^*}^2\qquad(x\in X).
 \label{eq:recourse-growth}
\end{equation}
Such an $x^*$ is a global minimizer. Residual ties are allowed; if the
core projection of $X$ is a single point, $g=+\infty$.

\begin{lemma}[Projected-growth tail]
\label{lem:recourse-growth}
\begin{enumerate}[label=(\alph*)]
\item For $t>0$ the set $G_t$ of draws admitting \eqref{eq:recourse-growth}
is closed, and $g(\gamma)\ge t$ if and only if $\gamma\in G_t$. If
$g(\gamma)>0$, all global minimizers have the same core $a(\gamma)$, and
$f_\gamma(x)-f_\gamma^*\ge t\norm{v-a(\gamma)}^2$ on $X$ for every
finite $t\in(0,g(\gamma)]$.
\item $\Pr_{\rm cont}\{g<t\}\le kt/\sigma$ for every $t>0$.
\item There is a base-computable $C_g\le2^{P_D(L)}$ with
$\Pr_{\rm grid}\{g<t\}\le kt/\sigma+C_g/M$ for every $t>0$ and every
power of two $M$.
\end{enumerate}
\end{lemma}

\begin{proof}
(a) Let $\gamma_l\to\gamma$ with witnesses $x^*_l$. A subsequence of the
witnesses converges in the compact set $X$, and \eqref{eq:recourse-growth}
passes to the limit; so $G_t$ is closed. If $g(\gamma)\ge t$, then
$\gamma\in G_{t'}$ for all $t'<t$; letting $t'\uparrow t$ and passing to
a convergent subsequence of witnesses shows $\gamma\in G_t$; the
converse holds by definition. If \eqref{eq:recourse-growth} holds with
$t>0$ and witness $x^*$, then $f_\gamma(x^*)=f_\gamma^*$, and every global
minimizer $x$ has $0\ge t\norm{v-v^*}^2$, so its core is $v^*$. Hence
$v^*=a(\gamma)$, and \eqref{eq:recourse-growth} is the displayed
inequality. If $g(\gamma)>0$, every finite $t\in(0,g(\gamma)]$ has
$\gamma\in G_t$, which gives the last claim.

(b) Translate the core coordinates by $(\frac12,\ldots,\frac12)$; this
changes $f_\gamma$ by a constant and leaves \eqref{eq:recourse-growth}
unchanged, and afterwards every core lies in $[-\frac12,\frac12]^k$.
Write $c=-\gamma$ and define the finite convex function
\[
 \Psi_t(c)=\max_{x\in X}\bigl(c^{\mathsf T}v-f(x)+t\norm v^2\bigr).
\]
As in Lemma~\ref{lem:recourse-conjugate}(a), every subgradient lies in
$[-\frac12,\frac12]^k$, and the core of every maximizer at $c$ is a
subgradient. Let $c$ be a point where $\Psi_t$ is differentiable, with
gradient $a$. Then every maximizer $\bar x$ has core $a$, and for
$y=c+2ta$ and every $x\in X$,
\[
 f(x)-f(\bar x)\ge c^{\mathsf T}(v-a)+t(\norm v^2-\norm a^2)
 =y^{\mathsf T}(v-a)+t\norm{v-a}^2 .
\]
Hence the draw $-y$ satisfies \eqref{eq:recourse-growth} with witness
$\bar x$. At differentiability points put $\hat c=c+2t\nabla\Psi_t(c)$;
then $\norm{\hat c-c}_\infty\le t$ and $-\hat c\in G_t$.

Let $P(y)$ be the unique minimizer of
$\Psi_t(c)+\norm{c-y}^2/(4t)$. Its optimality condition is
$y-P(y)\in2t\,\partial\Psi_t(P(y))$. If $y,y'$ are two points, adding the
two monotonicity inequalities gives
$\norm{P(y)-P(y')}^2\le(y-y')^{\mathsf T}(P(y)-P(y'))$, so $P$ is
$1$-Lipschitz. If $c$ is a differentiability point, then $c$ satisfies
the optimality condition for $y=\hat c$, so $P(\hat c)=c$.

Let $Q=[-\sigma,\sigma]^k$, $t<\sigma$, $Q'=[-\sigma+t,\sigma-t]^k$, and
let $\widetilde G=\{c\in Q:-c\in G_t\}$, a closed set. For every
differentiability point $c\in Q'$ we have $\hat c\in\widetilde G$ and
$P(\hat c)=c$. By (F1), the differentiability points have full measure, so
$P(\widetilde G)$ contains $Q'$ up to a null set. By (F2),
$\operatorname{Leb}(\widetilde G)\ge\operatorname{Leb}^*(P(\widetilde G))
\ge(2\sigma-2t)^k$. The uniform law of $c$ equals that of $\gamma$, so
$\Pr_{\rm cont}\{g<t\}\le1-(1-t/\sigma)^k\le kt/\sigma$. For
$t\ge\sigma$ the bound is trivial.

(c) $G_t$ is defined by
$\exists x^*\,\forall x:\ x^*\in X\wedge\bigl[x\notin X\vee
f_\gamma(x)-f_\gamma(x^*)-t\norm{v-v^*}^2\ge0\bigr]$, a formula
\eqref{eq:recourse-two-block} with two blocks of $n$ variables, at most
$2L+2$ atoms, and degree at most $\max\{D,2\}$. Its complement is
$\{g<t\}$. Lemmas~\ref{lem:recourse-sections}
and~\ref{lem:recourse-replacement} and part~(b) give the claim with
$C_g=2k(2H^2+1)$.
\end{proof}

The proof refines the classical observation that a linear tilt in
general position has a unique maximizer because convex conjugates are
differentiable almost everywhere. Here the conjugate includes the
growth term $t\norm v^2$, and the proximal map turns differentiability
into a quantitative bound.

\subsection{Proof of Theorem~\ref{thm:recourse-core}}
\label{app:recourse-core-proof}

\paragraph{Base choices.}
Let $\Xi_0$ be the fallback factor of Lemma~\ref{lem:recourse-fallback}
and let $C_W,C_g$ be the constants of
Lemmas~\ref{lem:recourse-count-formula} and~\ref{lem:recourse-growth}.
Choose, before sampling,
\begin{equation}
 \begin{gathered}
 \Xi\ge\max\{2,\Xi_0\}\ \text{with}\ \log\Xi\le P_D(L),\qquad
 g_0=\frac\sigma{2k\Xi},\qquad
 \Theta=4k+\frac{k^2\kappa_+}{g_0},\qquad N=4^k\Xi,\\
 M=\text{the least power of two with }M\ge2\Xi\max\{C_W,C_g,1\}.
 \end{gathered}
 \label{eq:recourse-core-params}
\end{equation}
All these numbers have bit length polynomial in $L$; in particular
$\log M\le P_D(L)$ and $b\le P_D(L)$. With $\beta=C_W/M\le1/(2\Xi)$,
Lemmas~\ref{lem:recourse-count-formula} and~\ref{lem:recourse-growth}
give
\begin{equation}
 \Pr\{W>s\}\le\frac{a_k}s+\beta\ \ (s\ge1),\qquad
 \Pr\{g<g_0\}\le\frac{kg_0}\sigma+\frac{C_g}M\le\frac1\Xi .
 \label{eq:recourse-core-tails}
\end{equation}
Neither the grid nor any other choice depends on $q$.

\paragraph{The evaluator.}
For a query $q$, let $J$ be the least integer $j\ge0$ with
$2e_{2^{-j}}\le2^{-q}$ and $\Theta2^{-j}\le2^{-q}$; then
$J\le q+P_D(L)$. Run the refinement with final level $J$ and cap $N$.
If it stops with \textsc{cap}, use the fallback. Otherwise let $H_J$ be
the coordinate hull of the cells retained at level $J$, that is, the
smallest box containing them, with widths $w_1,\ldots,w_k$. If
$\sum_iw_i^2\le2^{-2q}$, return $x_q=x_{\rm inc}$ and
$\ell_q=U-2e_{2^{-J}}$. Otherwise use the fallback. The \emph{fallback}
returns the output of Lemma~\ref{lem:recourse-fallback} for this draw
and this $q$.

\paragraph{Correctness on every draw.}
On the fallback branch, Lemma~\ref{lem:recourse-fallback} gives
\eqref{eq:recourse-core-output}. On the ordinary branch,
Lemma~\ref{lem:recourse-refinement}(a) gives
$\ell_q\le f_\gamma^*\le f_\gamma(x_q)=\ell_q+2e_{2^{-J}}\le\ell_q+2^{-q}$.
By Lemma~\ref{lem:recourse-refinement}(b) and~(c), the box $H_J$
contains $a(\gamma)$ and the core $v_q$ of $x_{\rm inc}$, so
$\norm{v_q-a(\gamma)}^2\le\sum_iw_i^2\le2^{-2q}$. Neither statement
uses uniqueness, growth, or probability. On the ordinary branch, $x_q$
is a witness that the cell oracle produced at a level at most $J$,
possibly an earlier one, and $\ell_q$ is computed from its value; both
have bit length $P_D(L+q)$.

\paragraph{When the ordinary branch fails.}
If the cap is reached before level $j$, then more than $2^k\Xi$ cells
were retained at level $j-1$, so $W>\Xi$ by
Lemma~\ref{lem:recourse-cell-count}. Suppose next that
$g(\gamma)\ge g_0$. By Lemma~\ref{lem:recourse-growth}(a) the optimal
core $a$ is unique and $f_\gamma(x)-f_\gamma^*\ge g_0\norm{v-a}^2$. With
$h=2^{-J}$, every retained cell $C$ has
$f_\gamma(\omega_C)-f_\gamma^*\le4e_h$ by
Lemma~\ref{lem:recourse-refinement}(d), hence
$\norm{v(\omega_C)-a}\le h\sqrt{\kappa_+k/(2g_0)}$, and every point of
$C$ is within $h$ of $v(\omega_C)$ in each coordinate. Therefore
$w_i\le2h(1+\sqrt{\kappa_+k/(2g_0)})$ for all $i$, and, using
$\sqrt k\le k$ and $\sqrt y\le1+y$,
\[
 \Bigl(\sum_iw_i^2\Bigr)^{1/2}\le2\sqrt kh\Bigl(1+\sqrt{\frac{\kappa_+k}{2g_0}}\Bigr)
 \le h\Bigl(4k+\frac{k^2\kappa_+}{g_0}\Bigr)=\Theta h\le2^{-q}.
\]
So the hull test passes. The fallback is used only on the event
$\{W>\Xi\}\cup\{g<g_0\}$, which does not depend on $q$.

\paragraph{Work.}
Without the cap, each level $j\le J$ generates at most
$4^k\min\{W,\Xi\}$ cells: at most $4^kW$ by
Lemma~\ref{lem:recourse-cell-count}, and at most $N=4^k\Xi$ because the
cap was not reached. Each cell costs $2^kP_D(L+b+j)$ by
Lemma~\ref{lem:recourse-cell-oracle}, and the hull is computed in one
pass. With $J\le q+P_D(L)$ and $b\le P_D(L)$, the ordinary work, and
the work spent before a fallback, is at most
$8^k\min\{W,\Xi\}P_D(L+q)$. The fallback costs at most
$\Xi P_D(L+q)$. Hence every query costs at most $\Omega(\gamma)P_D(L+q)$
with
\[
 \Omega=8^k\min\{W,\Xi\}+\Xi\,\mathbf1\{W>\Xi\ \text{or}\ g<g_0\}.
\]
By \eqref{eq:recourse-core-tails},
$\E\min\{W,\Xi\}=\int_0^\Xi\Pr\{\min\{W,\Xi\}>s\}\,ds
\le1+\int_1^\Xi(a_k/s+\beta)\,ds\le1+a_k\ln\Xi+\beta\Xi$ and
$\Xi\Pr\{W>\Xi\}\le a_k+\beta\Xi$. With $\beta\Xi\le\frac12$ and
$\Xi\Pr\{g<g_0\}\le1$,
\[
 \E\,\Omega\le8^k(2+a_k\ln\Xi)+a_k+2 .
\]
Since $1+\kappa_+/\sigma\le2(1+\kappa/\sigma)$, we have
$a_k\le(360k^3)^k(1+\kappa/\sigma)^k$, and $\ln\Xi\le P_D(L)$. This
proves $\E\,\Omega\le k^{O(k)}(1+\kappa/\sigma)^kP_D(L)$.
\qed

\begin{remark}[Low core dimension]
\label{rem:recourse-low-dim}
For $k\le2$ a direct count suffices. If $g>0$, the retained witnesses
lie within $h\sqrt{\kappa_+k/(2g)}$ of $a$, so at most
$4^k(3+\sqrt{2k\kappa_+/g})^k\le512\max\{1,\kappa_+/g\}^{k/2}$ cells are
generated per level. With $Z=\max\{1,\kappa_+/g\}^{k/2}$,
Lemma~\ref{lem:recourse-growth} gives
$\Pr\{Z>s\}\le k\kappa_+s^{-2/k}/\sigma+C_g/M$. Its truncated moment is
bounded for $k=1$ and logarithmic in the cap for $k=2$, which yields the
bound $(1+\kappa/\sigma)P_D(L+q)$. For $k\ge3$ the truncated moment grows
like a positive power of the cap, and the fallback term does as well;
this is why Lemmas~\ref{lem:recourse-conjugate}--\ref{lem:recourse-count-formula}
are needed.
\end{remark}

\subsection{Proof of Corollary~\ref{cor:recourse-joint}}
\label{app:recourse-joint}

Let $f$ be convex on $X$. Choose $\Xi\ge\max\{2,\Xi_0\}$, $g_0=\sigma/(2k\Xi)$,
and $M$ the least power of two with $M\ge2\Xi C_g$, so that
$\Pr\{g<g_0\}\le1/\Xi$ by Lemma~\ref{lem:recourse-growth}.

For a query $q$ put $\epsilon=2^{-q}$,
$\delta=\min\{\epsilon,g_0\epsilon^2/(128k)\}$, and
$\zeta=\epsilon/(8k)$. First apply \eqref{eq:recourse-cv} to $f_\gamma$
on $X$ with accuracy $\delta$, obtaining $y\in X$ and $\ell_0$ with
$\ell_0\le f_\gamma^*\le U:=f_\gamma(y)\le\ell_0+\delta$.

For a coordinate $i$ and a dyadic rational $\theta$, an \emph{upper
probe} applies \eqref{eq:recourse-cv} to $f_\gamma$ on the polyhedron
$X\cap\{v_i\ge\theta\}$ with accuracy $\delta$. It \emph{succeeds} if
this polyhedron is empty or the returned lower bound exceeds $U$;
otherwise it fails and returns a witness $x_\theta\in X$ with
$v_i(x_\theta)\ge\theta$ and $f_\gamma(x_\theta)\le U+\delta$. If an
upper probe succeeds at $\theta$, then every point of $X$ with
$v_i\ge\theta$ has value greater than $U$; hence every global minimizer
and $y$ itself have $v_i<\theta$.

The upper probe fails at $\theta=0$, because the lower bound is at most
$f_\gamma^*\le U$, and it succeeds at $\theta=2$. Binary search
maintains a failing $\mathrm{lo}$ with its witness and a succeeding
$\mathrm{hi}$, halving $\mathrm{hi}-\mathrm{lo}$ until it is at most
$\zeta$; set $u_i=\mathrm{hi}$. Lower probes on
$X\cap\{v_i\le\theta\}$, started from a succeeding $\theta=-1$ and a
failing $\theta=1$, give $l_i$ in the same way. No monotonicity of the
probe outcomes is needed. The box $\prod_i[l_i,u_i]$ contains the core
of $y$ and every optimal core. If $\sum_i(u_i-l_i)^2\le\epsilon^2$,
return $y$ and $\ell_0$; then \eqref{eq:recourse-core-output} holds.
Otherwise return the fallback of Lemma~\ref{lem:recourse-fallback}.

Suppose $g(\gamma)\ge g_0$. The witness at the final $\mathrm{lo}$ of
the upper search has value at most $U+\delta\le f_\gamma^*+2\delta$, so
its core is within $\sqrt{2\delta/g_0}\le\epsilon/(8\sqrt k)$ of $a$.
Hence $u_i\le\mathrm{lo}+\zeta\le a_i+\epsilon/(8\sqrt k)+\zeta$, and
symmetrically $l_i\ge a_i-\epsilon/(8\sqrt k)-\zeta$. Then
$u_i-l_i\le\epsilon/(4\sqrt k)+2\zeta\le\epsilon/(2\sqrt k)$, and the
test passes. The fallback is used only on $\{g<g_0\}$.

The ordinary branch makes $1+O(k(q+\log k))$ convex value calls on
polyhedra of bit length $P_D(L)+O(q+\log k)$, with tolerance
$\langle\delta\rangle\le2q+P_D(L)$. Since $k\le L$, its work is at most
$P_D(L+q)$, with a polynomial independent of $k$ and of the numerical
values of $\sigma$ and $\kappa$. With $\Omega=1+\Xi\mathbf1\{g<g_0\}$,
every query costs at most $\Omega(\gamma)P_D(L+q)$, and
$\E\,\Omega\le2$.
\qed

\subsection{Proof of Corollary~\ref{cor:recourse-qp}}
\label{app:recourse-qp}

\begin{lemma}[Height of a lexicographic quadratic optimizer]
\label{lem:recourse-qp-height}
Let $X=\{x\in\R^n:Ex\le e\}$ be a nonempty bounded polytope and
$q(x)=\frac12x^{\mathsf T}Ax+c^{\mathsf T}x+c_0$, with rational data and
$A=A^{\mathsf T}$, so that $\nabla q(x)=Ax+c$. Let
$\Delta_0$ be a positive integer making $\Delta_0A$, $\Delta_0c$,
$\Delta_0E$, and $\Delta_0e$ integral, and let $C\ge1$ bound the
absolute values of their entries. Then the lexicographically least
global minimizer $x^*$ of $q$ on $X$, for any fixed order of the
coordinates, has a common denominator at most
$U_Q=n!\,(n\cdot n!\,C^{n+1})^n$.
\end{lemma}

\begin{proof}
Let $I$ be the set of rows active at $x^*$, and let $R$ be a maximal
linearly independent subset of those rows. The smallest face of $X$
containing $x^*$ has affine hull $\{x:Rx=e_R\}$, and $x^*$ lies in its
relative interior, so $x^*\pm\epsilon w$ is feasible for small
$\epsilon$ and every $w\in\ker R$. Thus $w^{\mathsf T}(Ax^*+c)=0$ for
$w\in\ker R$. Let $N$ be an integer matrix whose columns span $\ker R$,
and put
\[
 S=\{x\in X:\ Rx=e_R,\ N^{\mathsf T}(Ax+c)=0\}\ni x^* .
\]
For $x,y\in S$ and $w=x-y\in\ker R$ we have
$w^{\mathsf T}(Ay+c)=0=w^{\mathsf T}(Ax+c)$, hence $w^{\mathsf T}Aw=0$ and
$q(x)-q(y)=w^{\mathsf T}(Ay+c)+\frac12w^{\mathsf T}Aw=0$. Thus every point
of $S$ is a global minimizer, and $x^*$ is the lexicographically least
point of $S$. A lexicographically least point of a polytope is a
vertex: one endpoint of any segment through it inside the polytope is
lexicographically smaller.

Choose $N$ from minors: if $\Delta_0R$ has rank $\rho$, fix a
nonsingular $\rho\times\rho$ submatrix, set one free coordinate equal to
its determinant and the other free coordinates to zero, and solve for
the pivot coordinates by the adjugate. The entries of $N$ are at most
$\rho!C^\rho\le n!C^n$. The equations $N^{\mathsf T}(\Delta_0Ax+\Delta_0c)=0$
then have integer coefficients bounded by $C_1=n\cdot n!C^{n+1}$, as do
the scaled rows of $E$. The vertex $x^*$ of $S$ is the unique solution
of $n$ linearly independent equations among these, and Cramer's rule
writes it with the common denominator $\abs{\det}\le n!C_1^n=U_Q$.
\end{proof}

\begin{proof}[Proof of Corollary~\ref{cor:recourse-qp}]
Order the core coordinates first; then the core of the lexicographically
least global minimizer is $a(\gamma)$. Write
$f_\gamma(x)=\frac12x^{\mathsf T}Ax+c^{\mathsf T}x+c_0$, where $A$ is the
constant symmetric Hessian of $f$ and only $c$ depends on $\gamma$. For
the grid $G_M$, every sampled coefficient has denominator dividing
$(M-1)$ times the denominator of $\sigma$, and its absolute value is at
most $\sigma$. Hence one base-computable $\Delta_0$ and $C$ serve all
draws in Lemma~\ref{lem:recourse-qp-height}, and $\log U_Q$ is
polynomial in $L$. Query Theorem~\ref{thm:recourse-core}, or
Corollary~\ref{cor:recourse-joint} for type~(P) with $\alpha=0$, at
precision $t$ with $2^{-t}\le1/(4U_Q^2)$, so $t\le P(L)$. Each interval
$[v_{t,i}-2^{-t},v_{t,i}+2^{-t}]$ contains $a_i$, a rational with
denominator at most $U_Q$; distinct such rationals differ by at least
$1/U_Q^2$, which exceeds the width $2^{1-t}$. The midpoint is within
$1/(4U_Q^2)<1/(2U_Q^2)$ of $a_i$, so by Legendre's theorem $a_i$ is a
convergent of its continued fraction; enumerating the convergents with
denominator at most $U_Q$ and keeping the one in the interval recovers
$a_i$ in polynomial time \citep{Schrijver1986}. Then
$X_a=\{z:(a,z)\in X\}$ is a nonempty bounded rational polyhedron on
which the quadratic $f_\gamma(a,\cdot)$ is convex, by residual convexity
in type~(B) and by Lemma~\ref{lem:recourse-selected-point}(b) in
type~(P). In type~(P) the set $X_a$ can be lower dimensional, and
convexity on $X_a$ does not make the constant Hessian of
$f_\gamma(a,\cdot)$ positive semidefinite on $\R^m$; for example, $-z^2$
is convex on $\{0\}$. We therefore first pass to the affine hull of
$X_a$. Lemma~\ref{lem:recourse-relative}, applied to $X_a$, gives a
rational point $z_0$, a rational matrix $B'$ of full column rank $m'$,
and the polytope $\mathcal R'=\{w:z_0+B'w\in X_a\}$, which contains a
ball about the origin and satisfies $X_a=z_0+B'\mathcal R'$. If $m'=0$,
put $z^*=z_0$. Otherwise the rational quadratic
$w\mapsto f_\gamma(a,z_0+B'w)$ is convex on $\mathcal R'$, which has
nonempty interior, so its constant Hessian $B'^{\mathsf T}A_{zz}B'$ is
positive semidefinite, where $A_{zz}$ is the residual block of $A$.
Solve this convex quadratic program over $\mathcal R'$ exactly
\citep{KozlovTarasovKhachiyan1980}, obtaining a rational optimal $w^*$,
and put $z^*=z_0+B'w^*$. Since $a$ is an optimal core,
$\min_{X_a}f_\gamma(a,\cdot)=f_\gamma^*$, so $(a,z^*)$ is a global
minimizer with value $f_\gamma^*=f_\gamma(a,z^*)$. The core query has
precision polynomial in $L$, and its output, even on the fallback
branch, is a rational of polynomial length; reconstruction, the linear
programs of Lemma~\ref{lem:recourse-relative}, and the exact quadratic
program cost polynomial time. The expected work is that of one query.
\end{proof}

\subsection{Error bounds and regularized selection}
\label{app:recourse-analytic}

The following lemmas are the forms of the shared error-bound and
selector tools (Lemma~\ref{lem:integer-hoffman},
Lemma~\ref{lem:selector}, and Lemmas~\ref{lem:points-cubic-transverse}
and~\ref{lem:points-cubic-slice}) needed here, with real coefficients, supplied
minor margins, and an inexact objective. We include their short proofs.

\begin{lemma}[Hoffman bound with a minor margin]
\label{lem:recourse-hoffman}
Let $B$ and $T$ be real matrices with $s$ columns and entries in
$[-C,C]$, $C\ge1$, and suppose every nonzero square minor of the stacked
matrix $\bigl[\begin{smallmatrix}B\\T\end{smallmatrix}\bigr]$ has
absolute value at least $\mu\in(0,1]$. Let $Q=\{x:Bx\le\beta\}$ and
$Z=\{x\in Q:Tx=t\}\ne\emptyset$, with arbitrary real $\beta$ and $t$.
Then $\dist(x,Z)\le(sC)^{s-1}\mu^{-1}\norm{Tx-t}$ for every $x\in Q$.
For integer matrices one may take $\mu=1$.
\end{lemma}

\begin{proof}
Let $w$ be the projection of $x$ onto the closed convex set $Z$; we may
assume $x\ne w$. Since $Z$ is polyhedral, $x-w=B_I^{\mathsf T}\lambda
+T^{\mathsf T}\nu$ with $\lambda\ge0$ and $I$ the rows of $B$ active at
$w$. Replace the rows of $T$ by a basis of its row space, and then
remove rows of $B_I$ until all remaining rows are linearly independent:
if a nontrivial combination of the remaining rows vanishes, it involves
some row of $B_I$, and subtracting a suitable multiple keeps all
$\lambda$ coefficients nonnegative while making one zero. We obtain an
independent row matrix $R$ with $r\le s$ rows, each a row of $B_I$ or
of $T$, and $x-w=R^{\mathsf T}\omega$ with nonnegative coefficients on
the rows from $B_I$. By the Cauchy--Binet formula, $\det(RR^{\mathsf T})$
is the sum of the squares of the $r\times r$ minors of $R$; at least
one is nonzero, so $\det(RR^{\mathsf T})\ge\mu^2$. Every eigenvalue of
$RR^{\mathsf T}$ is at most $\norm R_F^2\le s^2C^2$, so its least
eigenvalue is at least $\mu^2(sC)^{-2(r-1)}$, and
$\norm\omega\le(sC)^{s-1}\mu^{-1}\norm{x-w}$. Finally
$\norm{x-w}^2=\omega^{\mathsf T}R(x-w)$. A row of $B_I$ contributes
$\omega_i(B_ix-\beta_i)\le0$, since $B_iw=\beta_i$ and $x\in Q$; a row
of $T$ contributes $\omega_i(T_ix-t_i)$, since $T_iw=t_i$. Hence
$\norm{x-w}^2\le\norm\omega\norm{Tx-t}$, which gives the claim.
\end{proof}

\begin{lemma}[Cubic transverse estimate and slice]
\label{lem:recourse-transverse}
Let $\varphi$ be a real cubic polynomial on $\R^{n_0}$ and $Y\subseteq\R^{n_0}$
a convex set on which $\nabla^2\varphi\succeq0$. Let $c\in Y$,
$\rho\ge\sup_{x\in Y}\norm{x-c}$, and $M\ge\sup_Y\norm{\nabla^2\varphi}$.
Let $y$ minimize $\varphi$ on $Y$, let $x\in Y$, $d=x-y$, and
$E=\varphi(x)-\varphi(y)$.
\begin{enumerate}[label=(\alph*)]
\item $\bigl(d^{\mathsf T}\nabla^2\varphi(c)d\bigr)^2\le108M\rho^2E$.
\item Let $K$ be a subspace contained in $\ker\nabla^2\varphi(\xi)$ for
every $\xi\in Y$, $\Pi$ the orthogonal projection onto $K^\perp$, and
$g=\nabla\varphi(c)$. For $x',y'\in Y$,
$\abs{\varphi(x')-\varphi(y')-g^{\mathsf T}(x'-y')}\le M\rho\norm{\Pi(x'-y')}$.
In particular $\abs{g^{\mathsf T}d}\le E+M\rho\norm{\Pi d}$.
\item If moreover $\ker\nabla^2\varphi(c)=K$, then
$\argmin_Y\varphi=\{x\in Y:\nabla^2\varphi(c)(x-y)=0,\ g^{\mathsf T}(x-y)=0\}$.
\end{enumerate}
\end{lemma}

\begin{proof}
(a) Put $a_0=d^{\mathsf T}\nabla^2\varphi(y)d\ge0$,
$b_0=d^{\mathsf T}\nabla^2\varphi(x)d\ge0$, and $u=a_0+b_0$. Since
$\varphi$ is cubic, $\nabla^2\varphi(x)-\nabla^2\varphi(y)=\mathcal T[d]$
for the constant symmetric third-derivative tensor $\mathcal T$, and
exact Taylor expansion gives
$E=\nabla\varphi(y)^{\mathsf T}d+\frac12a_0+\frac16(b_0-a_0)\ge u/6$,
because $\nabla\varphi(y)^{\mathsf T}d\ge0$ by first-order optimality on
the convex set $Y$. Symmetry of $\mathcal T$ gives
$d^{\mathsf T}\nabla^2\varphi(c)d=a_0+(c-y)^{\mathsf T}
(\nabla^2\varphi(x)-\nabla^2\varphi(y))d$. For positive semidefinite
$P$, $\norm{Pd}\le\sqrt{\norm P\,d^{\mathsf T}Pd}$; with
$\norm{c-y}\le\rho$ this gives
$d^{\mathsf T}\nabla^2\varphi(c)d\le u+\rho\sqrt M(\sqrt{a_0}+\sqrt{b_0})
\le u+\rho\sqrt{2Mu}$. Since $\norm d\le2\rho$, $u\le8M\rho^2$, so
$u\le2\rho\sqrt{2Mu}$. As $c\in Y$, the left side is nonnegative, and
$(d^{\mathsf T}\nabla^2\varphi(c)d)^2\le18M\rho^2u\le108M\rho^2E$.

(b) For $\xi\in Y$, $\nabla\varphi(\xi)-g=\int_0^1\nabla^2\varphi(c+\theta(\xi-c))
(\xi-c)\,d\theta$. Each Hessian on $Y$ is symmetric with $K$ in its
kernel, so its range lies in $K^\perp$; hence $\nabla\varphi(\xi)-g\in
K^\perp$ and $\norm{\nabla\varphi(\xi)-g}\le M\rho$. Integrating along
the segment from $y'$ to $x'$ gives
$\varphi(x')-\varphi(y')-g^{\mathsf T}(x'-y')=\int_0^1
(\nabla\varphi(y'+\theta(x'-y'))-g)^{\mathsf T}\Pi(x'-y')\,d\theta$.

(c) If $x$ is a minimizer, then $E=0$, so (a) and positive
semidefiniteness give $\nabla^2\varphi(c)d=0$; thus $d\in K$, $\Pi d=0$,
and (b) gives $g^{\mathsf T}d=0$. Conversely, the two equations give
$\Pi d=0$ and then $\varphi(x)-\varphi(y)=g^{\mathsf T}d=0$ by (b).
\end{proof}

\begin{lemma}[Gradient rows under global convexity]
\label{lem:recourse-bregman}
Let $D\ge2$ and let $\varphi$ be a real polynomial on $\R^{n_0}$ of degree
at most $D$ that is convex on all of $\R^{n_0}$. Let $Y\subseteq\R^{n_0}$ be
convex, $y$ a minimizer of $\varphi$ on $Y$ with $\norm y_\infty\le\rho$,
$\rho\ge1$, and $x\in Y$ with $E=\varphi(x)-\varphi(y)\in[0,1]$. Let
$V\ge\sup\{\abs{\varphi(\xi)}:\norm\xi_\infty\le\rho+2\}$ and
$G\ge\sup\{\norm{\nabla\varphi(\xi)}_1:\norm\xi_\infty\le\rho\}$. Then
every coefficient of the polynomial $\xi\mapsto\nabla\varphi(\xi)^{\mathsf T}(x-y)$
has absolute value at most
\[
 C_\varphi E^{1/D},\qquad
 C_\varphi=B_{D-1}^{n_0}\bigl(1+J_D(V+G(1+\rho)+C_D3^D)\bigr),
\]
where $C_D=(D+1)D^D$, $J_D=(D+1)D^{D+1}$, and $B_r=(r+1)2^rr^r$.
\end{lemma}

\begin{proof}
Let $d=x-y$ and let
$\beta(u)=\varphi(y+u)-\varphi(y)-\nabla\varphi(y)^{\mathsf T}u$, which is
convex and nonnegative on $\R^{n_0}$ with $\beta(0)=0$. Since
$\nabla\varphi(y)^{\mathsf T}d\ge0$ and $\beta(d)\ge0$, we have
$0\le\nabla\varphi(y)^{\mathsf T}d\le E$ and $\beta(d)\le E$; convexity
gives $0\le\beta(\theta d)\le E$ for $\theta\in[0,1]$. Two
interpolation facts are used. If a univariate polynomial of degree at
most $D$ is bounded by $E$ at the nodes $j/D$, $0\le j\le D$, then it is
bounded by $C_DE(1+\abs\theta)^D$ at every real $\theta$, because each
Lagrange basis polynomial has denominator at least $D^{-D}$ and
numerator at most $(1+\abs\theta)^D$. If a polynomial $\pi$ of degree
at most $D$ is bounded by $K$ on $[0,\bar\theta]$, then
$\abs{\pi'(0)}\le J_DK/\bar\theta$, by interpolating at the nodes
$j\bar\theta/D$.

Fix $\xi\in[0,1]^{n_0}$ and let $\pi(\theta)=\beta(\xi-y+\theta d)$. By
midpoint convexity,
$0\le\pi(\theta)\le\frac12\beta(2(\xi-y))+\frac12\beta(2\theta d)$. Here
$\beta(2(\xi-y))\le2V+2G(1+\rho)$, because $\norm{2\xi-y}_\infty\le\rho+2$
and $\norm{\xi-y}_\infty\le1+\rho$, and
$\beta(2\theta d)\le C_DE(1+2\abs\theta)^D$. If $E>0$, then on
$[0,E^{-1/D}]$ we have $E(1+2\theta)^D\le(E^{1/D}+2)^D\le3^D$, so
$\abs\pi\le K:=V+G(1+\rho)+C_D3^D/2$ there and
$\abs{\pi'(0)}\le J_DKE^{1/D}$. If $E=0$, then $\beta(\theta d)=0$ on
$[0,1]$ and hence for all $\theta$, so $\pi$ is bounded on $\R$ and
constant, and $\pi'(0)=0$. Since
$\pi'(0)=(\nabla\varphi(\xi)-\nabla\varphi(y))^{\mathsf T}d$, in both cases
$\abs{\nabla\varphi(\xi)^{\mathsf T}d}\le(1+J_DK)E^{1/D}$ on $[0,1]^{n_0}$.

Finally, a polynomial of degree at most $r$ in each of $n_0$ variables
that is bounded by $K'$ on $[0,1]^{n_0}$ has all coefficients bounded by
$B_r^{n_0}K'$: interpolate on the grid $\{0,1/r,\ldots,1\}^{n_0}$, and note that
each univariate Lagrange basis polynomial on these nodes has
coefficients at most $2^rr^r$. Apply this with $r=D-1$.
\end{proof}

\begin{lemma}[Regularized selection with an inexact objective]
\label{lem:recourse-selection}
Let $Y\subseteq\R^{n_0}$ be compact and convex, $\varphi$ convex and
continuous on $Y$, $S=\argmin_Y\varphi$, $\varphi^*=\min_Y\varphi$, and $p$
the least-norm point of $S$. Let $R_x\ge\max\{1,\sup_Y\norm x\}$,
$\Gamma\ge1$, and $D_e\ge2$, and suppose
$\dist(x,S)\le\Gamma(\varphi(x)-\varphi^*)^{1/D_e}$ whenever $x\in Y$ and
$\varphi(x)-\varphi^*\le1$. For $0<e\le1$ put
\[
 \tau=\frac{e^{2D_e-2}}{2^{3D_e-2}R_x^{D_e}\Gamma^{D_e}} .
\]
Then every $\hat x\in Y$ with
$\varphi(\hat x)+\tau\norm{\hat x}^2\le\min_Y(\varphi+\tau\norm\cdot^2)+\tau e^2/4$
satisfies $\norm{\hat x-p}\le e$.
\end{lemma}

\begin{proof}
Let $x_\tau$ minimize $\varphi+\tau\norm\cdot^2$ on $Y$. Comparison with
$p$ gives $\norm{x_\tau}\le\norm p$ and
$\varphi(x_\tau)-\varphi^*\le\tau R_x^2\le1$. Let $s$ be the point of $S$
nearest to $x_\tau$ and $\epsilon_d=\norm{x_\tau-s}$. Comparison with $s$
gives $\varphi(x_\tau)-\varphi^*\le\tau(\norm s^2-\norm{x_\tau}^2)\le
2R_x\tau\epsilon_d$, and the error bound gives
$\epsilon_d^{D_e}\le\Gamma^{D_e}2R_x\tau\epsilon_d$, so
$\epsilon_d\le(2R_x\tau\Gamma^{D_e})^{1/(D_e-1)}=e^2/(8R_x)$. Since $p$
is the projection of the origin onto $S$, $p^{\mathsf T}(s-p)\ge0$, and
\[
 \norm{x_\tau-p}^2\le2\norm p^2-2p^{\mathsf T}x_\tau
 =2p^{\mathsf T}(p-s)+2p^{\mathsf T}(s-x_\tau)\le2R_x\epsilon_d\le e^2/4 .
\]
The function $\varphi+\tau\norm\cdot^2$ is $2\tau$-strongly convex, so
$\tau\norm{\hat x-x_\tau}^2\le\tau e^2/4$. The triangle inequality gives
$\norm{\hat x-p}\le e$.
\end{proof}

\begin{lemma}[Relative coordinates]
\label{lem:recourse-relative}
For a nonempty bounded rational polytope $X=\{x\in\R^n:Ex\le e\}$ one
computes in polynomial time an integer $n'\ge0$, a rational $x_0\in X$,
a rational $n\times n'$ matrix $B$ of rank $n'$, and rationals $r>0$ and
$R_0\ge1$, all of polynomial bit length, such that
$\mathcal R=\{w:x_0+Bw\in X\}$ satisfies $X=x_0+B\mathcal R$ and
$B(0,r)\subseteq\mathcal R\subseteq B(0,R_0)$.
\end{lemma}

\begin{proof}
For each row $i$ solve the linear program $\max\{e_i-E_ix:x\in X\}$. Let
$I_0$ be the rows whose maximum slack is zero; then
$\operatorname{aff}X=\{x:E_{I_0}x=e_{I_0}\}$. If every row lies in
$I_0$, $X$ is a point and $n'=0$. Otherwise average optimal solutions
for the rows outside $I_0$ to get $\bar x\in X$ with positive slack in
all of them. Let the columns of $B$ be a rational basis of
$\ker E_{I_0}$. The linear program
$\max\{\varrho:\ E_i(\bar x+Bw)+\varrho\norm{E_iB}_1\le e_i\ (i\notin I_0),
\ 0\le\varrho\le1\}$ has a positive optimum $r$ at a rational point $w^*$;
put $x_0=\bar x+Bw^*$. Every $u$ with $\norm u_\infty\le r$ keeps
$x_0+Bu$ in $X$, so the closed, and hence the open, Euclidean ball of
radius $r$ lies in $\mathcal R$. Coordinate linear programs bound
$\abs{w_j}$ on $\mathcal R$ by rationals $m_j$. Put $R_0=1+\sum_jm_j$;
then $\norm w\le\norm w_1\le\sum_jm_j<R_0$ on $\mathcal R$, so
$\mathcal R\subseteq B(0,R_0)$.
\end{proof}

\subsection{Proof of Theorem~\ref{thm:recourse-convexified}}
\label{app:recourse-convexified}

Throughout this subsection the instance has type~(P), $\kappa=\alpha$,
and we assume $D\ge2$ (a degree bound can always be raised). Put
$G_\alpha=f+\frac\alpha2\norm v^2$, $\beta=\alpha+1$, and
$G_\beta=G_\alpha+\frac12\norm v^2$. For a draw $\gamma$ with selected
core $a$ define the surrogate
\[
 \psi_a(x)=f_\gamma(x)+\frac\beta2\norm{v-a}^2
 =G_\beta(x)+(\gamma-\beta a)^{\mathsf T}v+\frac\beta2\norm a^2,
 \qquad S_a=\{a\}\times S(\gamma).
\]

\begin{lemma}[The convex surrogate]
\label{lem:recourse-surrogate}
In case~(i) $\psi_a$ is convex on $X$; in case~(ii) it is convex on
$\R^n$. Moreover $\min_X\psi_a=f_\gamma^*$, $\argmin_X\psi_a=S_a$,
$\psi_a(x)-f_\gamma^*\ge\frac\beta2\norm{v-a}^2$ on $X$, and the
least-norm point of $S_a$ is $x^\sharp(\gamma)$.
\end{lemma}

\begin{proof}
$\psi_a$ equals $G_\alpha$ plus the convex function $\frac12\norm v^2$
plus an affine function. On $X$, $\psi_a(x)\ge f_\gamma(x)\ge f_\gamma^*$,
with equality exactly when $f_\gamma(x)=f_\gamma^*$ and $v=a$, that is,
on $S_a$. The least-norm point of $\{a\}\times S(\gamma)$ is $(a,p)$.
\end{proof}

\begin{lemma}[A uniform error bound with rational rows]
\label{lem:recourse-convexified-error}
Let $D_e=4$ in case~(i) and $D_e=D$ in case~(ii). There is a
base-computable rational $\Gamma\ge1$ of polynomial bit length such
that, for every $\gamma\in[-\sigma,\sigma]^k$ and every $x\in X$ with
$E=\psi_a(x)-f_\gamma^*\le1$,
\[
 \dist(x,S_a)\le\Gamma E^{1/D_e}.
\]
\end{lemma}

\begin{proof}
Put $N=k(\sigma+\beta)$, so that $\norm{\gamma-\beta a}\le N$. Fix
$y\in S_a$ and write $d=x-y$ and $d_v=v-a$ for its core part. By
Lemma~\ref{lem:recourse-surrogate}, $\norm{d_v}\le\sqrt{2E/\beta}\le
2E^{1/D_e}$ when $E\le1$.

\emph{Case~(i).} Take $x_0,B,r,R_0,\mathcal R,n'$ from
Lemma~\ref{lem:recourse-relative}. If $n'=0$, then $X$ is a point and
$\Gamma=1$ works. Otherwise let $B_v$ be the first $k$ rows of $B$, so
that the core of $x_0+Bw$ is $v_0+B_vw$. The function
$h(w)=G_\beta(x_0+Bw)$ is an explicit rational cubic; put
$H(w)=\nabla^2h(w)$, $H_0=H(0)$, $g_0=\nabla h(0)$, and $K=\ker H_0$. Since
$h$ is convex on the full-dimensional set $\mathcal R$, $H\succeq0$ on
$\mathcal R$. For $w\in\mathcal R$ the point $w'=-(r/R_0)w$ lies in
$B(0,r)\subseteq\mathcal R$, and $0$ is a convex combination of $w$ and
$w'$ with weights $\theta/(1+\theta)$ and $1/(1+\theta)$, $\theta=r/R_0$.
As $H$ is affine, $H_0\succeq\frac\theta{1+\theta}H(w)$, that is,
\[
 0\preceq H(w)\preceq\beta_0H_0\quad(w\in\mathcal R),\qquad
 \beta_0=1+R_0/r .
\]
Hence $K\subseteq\ker H(w)$ on $\mathcal R$. Let
$M_0=\max\{1,\max_i\sum_j\abs{(H_0)_{ij}}\}\ge\norm{H_0}$ and
$M=\beta_0M_0$. If $H_0\ne0$, let $\Delta_H$ be a positive integer with
$\Delta_HH_0$ integral, let $C_H\ge1$ bound its entries, and put
$\lambda_0=1/(\Delta_H(n'C_H)^{n'-1})$; the product of the positive
eigenvalues of $\Delta_HH_0$ is a sum of principal minors, hence a
positive integer, and each eigenvalue is at most $n'C_H$, so every
positive eigenvalue of $H_0$ is at least $\lambda_0$. Put
$C_\perp=\max\{1,108MR_0^2/\lambda_0^2\}$, or $C_\perp=1$ if $H_0=0$.

Let $w,u\in\mathcal R$ represent $x,y$, put $d'=w-u$, and let $\Pi$
project onto $K^\perp$. The reduced surrogate
$\tau_a(w)=\psi_a(x_0+Bw)=h(w)+(\gamma-\beta a)^{\mathsf T}(v_0+B_vw)+{\rm const}$
is a cubic, convex on $\mathcal R$, with Hessian $H$ and minimizer $u$,
and $E=\tau_a(w)-\tau_a(u)$. Lemma~\ref{lem:recourse-transverse}(a) with
$c=0$ and $\rho=R_0$ gives $(d'^{\mathsf T}H_0d')^2\le108MR_0^2E$; if
$H_0\ne0$, then $d'^{\mathsf T}H_0d'\ge\lambda_0\norm{\Pi d'}^2$, so
$\norm{\Pi d'}\le C_\perp E^{1/4}$, and if $H_0=0$ then $\Pi=0$.
Lemma~\ref{lem:recourse-transverse}(b), applied to $h$, gives
$\abs{h(w)-h(u)-g_0^{\mathsf T}d'}\le MR_0\norm{\Pi d'}$. Since
$E=h(w)-h(u)+(\gamma-\beta a)^{\mathsf T}B_vd'$ and $B_vd'=d_v$,
\[
 \abs{g_0^{\mathsf T}d'}\le E+N\norm{d_v}+MR_0\norm{\Pi d'} .
\]
Let $J=[H_0;g_0^{\mathsf T};B_v]$, a known rational matrix. For $E\le1$,
$\norm{Jd'}\le M_0\norm{\Pi d'}+\abs{g_0^{\mathsf T}d'}+\norm{d_v}\le
C_{\rm res}E^{1/4}$ with $C_{\rm res}=3+2N+(M_0+MR_0)C_\perp$.

Let $\mathcal R_a=\{w\in\mathcal R:x_0+Bw\in S_a\}$ be the reduced
optimal set. Then $\mathcal R_a=\{w\in\mathcal R:J(w-u)=0\}$. Indeed, if
$w$ is optimal, that is, $w\in\mathcal R_a$, then $E=0$ and the bounds
give $Jd'=0$. Conversely, $Jd'=0$ gives $d'\in K$, so
$h(w)-h(u)=g_0^{\mathsf T}d'=0$ by
Lemma~\ref{lem:recourse-transverse}(b), and $B_vd'=0$, so $E=0$ and $w$
is optimal. Clear
the denominators of $J$ and of the rows $EB$ describing $\mathcal R$ by
a positive integer $\Delta_J$, and let $C_J\ge1$ bound the entries of the
two integer matrices. Lemma~\ref{lem:recourse-hoffman} with $\mu=1$ gives
$\dist(w,\mathcal R_a)\le(n'C_J)^{n'-1}\Delta_J\norm{Jd'}$. Since
$x=x_0+Bw$, the claim holds with
$\Gamma=B_F(n'C_J)^{n'-1}\Delta_JC_{\rm res}$, where $B_F\ge\max\{1,\norm B\}$
is a rational bound, for example $B_F=\max\{1,\sum_{i,j}\abs{B_{ij}}\}$.

\emph{Case~(ii).} Let $\rho\ge1$ bound $\norm x_\infty$ on the supplied
box. Writing $f=\sum_\nu f_\nu x^\nu$, the numbers
\[
 V=\sum_\nu\abs{f_\nu}(\rho+2)^{\abs\nu}+N(\rho+2)+\tfrac\beta2k(\rho+3)^2,
 \qquad
 G=\sum_\nu\abs{f_\nu}\abs\nu\rho^{\abs\nu-1}+k\sigma+\beta k(\rho+1)
\]
bound $\abs{\psi_a}$ on $\{\norm\xi_\infty\le\rho+2\}$ and
$\norm{\nabla\psi_a}_1$ on $\{\norm\xi_\infty\le\rho\}$, uniformly in $a$
and $\gamma$. Lemma~\ref{lem:recourse-bregman}, applied to $\psi_a$ on
$Y=X$, bounds every coefficient of $\xi\mapsto\nabla\psi_a(\xi)^{\mathsf T}d$
by $C_\psi E^{1/D}$, with $C_\psi$ computed from $V$, $G$, $\rho$, $n$, $D$.
For a polynomial $\varphi$ of degree at most $D$ let $N_\varphi$ be the
matrix whose rows are indexed by the monomials $\xi^\mu$ of degree at
most $D-1$, with row $\mu$ equal to $((\mu_i+1)\varphi_{\mu+e_i})_{i\le n}$;
then $N_\varphi d$ is the coefficient vector of
$\nabla\varphi(\xi)^{\mathsf T}d$. Since $\psi_a-G_\beta$ is affine,
$N_{\psi_a}$ and $N_{G_\beta}$ differ only in the row $\mu=0$, by
$((\gamma-\beta a)^{\mathsf T},0)$. Let $J=[N_{G_\beta};\,(I_k\ 0)]$,
a known rational matrix with at most $(n+1)^{D-1}+k$ rows. The vector
$N_{\psi_a}d$ has at most $\binom{n+D-1}{D-1}\le(n+1)^{D-1}$ entries, each
of absolute value at most $C_\psi E^{1/D}$, so its norm is at most
$(n+1)^{\lceil(D-1)/2\rceil}C_\psi E^{1/D}$. The integer exponent makes
the constant $C_{\rm res}$ in the next display rational. Then, for
$E\le1$,
\[
 \norm{Jd}\le(n+1)^{\lceil(D-1)/2\rceil}C_\psi E^{1/D}+N\norm{d_v}+\norm{d_v}
 \le C_{\rm res}E^{1/D},\qquad
 C_{\rm res}=(n+1)^{\lceil(D-1)/2\rceil}C_\psi+2N+2 .
\]
The optimal set is $S_a=\{x\in X:J(x-y)=0\}$: if $E=0$ the bound gives
$Jd=0$; conversely $Jd=0$ gives $d_v=0$ and $N_{G_\beta}d=0$, hence
$N_{\psi_a}d=0$, so $\nabla\psi_a(\xi)^{\mathsf T}d=0$ for all $\xi$ and
$\psi_a(x)=\psi_a(y)$. Lemma~\ref{lem:recourse-hoffman}, applied after
clearing the denominators of $J$ and of the rows of $X$, gives the claim
with $\Gamma=(nC_J)^{n-1}\Delta_JC_{\rm res}$.

In both cases the constants depend only on the instance, not on
$\gamma$ or $a$, and their bit lengths are polynomial in $L$ for fixed
$D$.
\end{proof}

\begin{proof}[Proof of Theorem~\ref{thm:recourse-convexified}]
Use the grid and the core evaluator of Theorem~\ref{thm:recourse-core}
with $\kappa=\alpha$, or of Corollary~\ref{cor:recourse-joint} if
$\alpha=0$. For a query $q$ let $\epsilon=2^{-q}$, let $R_x\ge
\max\{1,\sup_X\norm x\}$ be a rational bound from the supplied box, and
put
\[
 \tau=\frac{\epsilon^{2D_e-2}}{2^{3D_e-2}R_x^{D_e}\Gamma^{D_e}},\qquad
 \eta_q=\frac{\tau\epsilon^2}8,\qquad
 \delta_q=\frac{\tau\epsilon^2}{16\beta k}.
\]
Query the core evaluator with accuracy $\delta_q$; it returns a rational
$b$ with $\norm{b-a}\le\delta_q$. Apply \eqref{eq:recourse-cv} on $X$ to
the explicit rational convex polynomial
$\Phi_b(x)=G_\beta(x)+(\gamma-\beta b)^{\mathsf T}v+\tau\norm x^2$ with
accuracy $\eta_q$, and return its point $\hat x\in X$. Let
$\Phi_a=\psi_a-\frac\beta2\norm a^2+\tau\norm x^2$. On $X$,
$\abs{\Phi_b-\Phi_a}=\beta\abs{(a-b)^{\mathsf T}v}\le\beta\sqrt k\,\delta_q$,
so
$\Phi_a(\hat x)-\min_X\Phi_a\le\eta_q+2\beta\sqrt k\,\delta_q\le\tau\epsilon^2/4$.
By Lemmas~\ref{lem:recourse-surrogate}, \ref{lem:recourse-convexified-error},
and~\ref{lem:recourse-selection} with $e=\epsilon$,
$\norm{\hat x-x^\sharp(\gamma)}\le\epsilon$. The core query has precision
$\log(1/\delta_q)\le(2D_e+2)q+P_D(L)$ and costs at most
$\Omega(\gamma)P_D(L+q)$. The convex solve has data of bit length
$P_D(L+q)$ and tolerance $\langle\eta_q\rangle\le(2D_e+2)q+P_D(L)$, so it
costs $P_D(L+q)$. The random
factor is that of the core evaluator, so its expectation has the stated
bound.
\end{proof}

The completion never reads an exact algebraic record of $a$: it uses
only the short rational output $b$. The surrogate is optimized on the
fixed domain $X$, so no continuity of residual optimizers in the core is
needed.

\begin{proof}[Proof of Corollary~\ref{cor:recourse-affine-power}]
Each $w_j(a_j^{\mathsf T}x-b_j)^{p_j}$ is convex because $p_j$ is even
and $w_j>0$, and $\frac12x^{\mathsf T}Ax$ is convex because
$A\succeq0$; so the representation is convex on $\R^n$. For fixed $D$,
expanding the right side produces at most $\binom{n+D}D$ monomials, so
the identity with $f+\frac\alpha2\norm v^2$ is checked in polynomial
time, and $A\succeq0$ is checked by exact symmetric Gaussian elimination.
\end{proof}

\subsection{Residual cubic fibers}
\label{app:recourse-fiber}

\begin{proof}[Proof of Proposition~\ref{prop:recourse-fiber}]
If $m=0$, the residual space consists of the empty vector. Then
$S_v=\{()\}$, all matrix and kernel assertions are vacuous, and part~(c)
holds with $\Gamma=1$. In the rest of the proof assume $m\ge1$.

(a) Since $f$ is cubic, $H$ is affine in $(v,z)$. For each
$v\in[0,1]^k$, $f(v,\cdot)$ is convex on $[0,1]^m$, so $H(v,z)\succeq0$
there, by continuity at the boundary. The reflection
$(v,z)\mapsto(2c_A-v,\mathbf1-z)$ maps $A\times[0,1]^m$ onto itself, and
affinity gives $H(v,z)+H(2c_A-v,\mathbf1-z)=2H_A$. Both terms are
positive semidefinite, so $0\preceq H(v,z)\preceq2H_A$, and
$K_A:=\ker H_A\subseteq\ker H(v,z)$. If the free coordinates of $v$ lie in
$[\delta,1-\delta]$ with $\delta<\frac12$, then
$w=(v-2\delta c_A)/(1-2\delta)\in A$ and $v=2\delta c_A+(1-2\delta)w$, so
$H(v,c_z)=2\delta H_A+(1-2\delta)H(w,c_z)\succeq2\delta H_A$; for
$\delta=\frac12$, $v=c_A$. Every $v$ in the relative interior of $A$ has
such a $\delta>0$, so $\ker H(v,c_z)\subseteq K_A$.

(b) Apply Lemma~\ref{lem:recourse-transverse}(c) to $\varphi=f(v,\cdot)$
on $Y=[0,1]^m$ with $c=c_z$, $\rho=\sqrt m/2$, and $K=K_A$; part~(a)
gives the kernel hypotheses. Since $\ker H(v,c_z)=K_A=\ker H_A$, the
condition $H(v,c_z)(z-y)=0$ is equivalent to $H_A(z-y)=0$. A square
submatrix of $J_A(v)$ contains the row $g_A(v)^{\mathsf T}$ at most once;
its determinant is constant if it does not, and otherwise a linear
combination, with constant coefficients, of the entries of $g_A(v)$,
which have degree at most two in $v$.

(c) Let $M_A=\max\{1,\max_j\sum_l\abs{(H_A)_{jl}}\}$ and $M_*=2M_A$, so
that $\norm{H(v,z)}\le M_*$ on $A\times[0,1]^m$. If $H_A\ne0$, let
$\lambda_0^A$ be the rational lower bound on its positive eigenvalues
from principal minors, as in the proof of
Lemma~\ref{lem:recourse-convexified-error}, and put
$C_\perp=\max\{1,27M_*m/(2\delta\lambda_0^A)^2\}$; if $H_A=0$ put
$C_\perp=1$. Let $C_B\ge1$ bound the absolute values of the entries of
$H_A$ and of $g_A(v')$ for all $v'\in[0,1]^k$, by coefficient sums. Fix
$z\in[0,1]^m$ with gap $E\le1$, $y\in S_v$, $d=z-y$, and $\Pi$ the
projection onto $K_A^\perp$. We first show
$\norm{\Pi d}\le C_\perp E^{1/4}$. If $H_A=0$, then $K_A=\R^m$ and
$\Pi=0$, so there is nothing to prove. If $H_A\ne0$, part~(a) gives
$d^{\mathsf T}H(v,c_z)d\ge2\delta\,d^{\mathsf T}H_Ad
\ge2\delta\lambda_0^A\norm{\Pi d}^2$, and
Lemma~\ref{lem:recourse-transverse}(a) with $\rho^2=m/4$ gives
$(d^{\mathsf T}H(v,c_z)d)^2\le27M_*mE$; together these give the claim.
Part~(b) of that lemma, with $\rho=\sqrt m/2\le m$, gives
$\abs{g_A(v)^{\mathsf T}d}\le E+M_*m\norm{\Pi d}$. Using $m$ in place of
$\sqrt m/2$ makes the constant $\Gamma$ defined next rational. With
$T=[H_A;g_A(v)^{\mathsf T}]$ and
$H_Ad=H_A\Pi d$,
$\norm{Td}\le\bigl(1+(M_A+M_*m)C_\perp\bigr)E^{1/4}$.
By part~(b), $S_v=\{z\in[0,1]^m:Tz=Ty\}$. Write
$[0,1]^m=\{z:z\le\mathbf1,\,-z\le0\}$. Up to sign, every nonzero square
minor of the stacked matrix of box rows and $T$ is a nonzero square
minor of $J_A(v)$, hence at least $\mu$ in absolute value.
Lemma~\ref{lem:recourse-hoffman} gives
\[
 \dist(z,S_v)\le\Gamma E^{1/4},\qquad
 \Gamma=(mC_B)^{m-1}\mu^{-1}\bigl(1+(M_A+M_*m)C_\perp\bigr).
\]
All quantities in this formula are rational and computed by arithmetic
on $f$, $A$, $\delta$, and $\mu$, so $\Gamma$ is computed in time, and
has bit length, polynomial in $L+\langle\delta\rangle+\langle\mu\rangle$.
\end{proof}

\begin{lemma}[Uniform heights of the minors]
\label{lem:recourse-minor-heights}
There are base-computable positive integers $\Delta$ and $B$, with
$\log\Delta+\log B$ polynomial in $L$, such that for every face $A$ and
every square submatrix of $J_A(v)$, the polynomial $\Delta\det(\cdot)$,
with the fixed coordinates of $A$ substituted, has integer coefficients
of absolute value at most $B$ in the free coordinates.
\end{lemma}

\begin{proof}
Let $d_0$ be the product of the denominators of the coefficients of
$f=\sum_\nu f_\nu x^\nu$, and $d=8d_0$. Entries of $H_A$ are values at
points with coordinates in $\{0,\frac12,1\}$ of second derivatives of
degree at most one, so $d$ clears them. Entries of $g_A(v)$ are
first $z$-derivatives, of degree at most two, evaluated at $z=c_z$; after
substituting the fixed coordinates of $A$, their coefficients are
cleared by $d$. Let $C_0=\max\{d,\lceil6d\sum_\nu\abs{f_\nu}\rceil\}$.
Then $C_0$ bounds the entries of $dH_A$, the entries $d$ of $d\,I_m$,
and the sum of the absolute coefficients of each entry of $d\,g_A$ on
every face. A minor of order $r'\le m$ is a sum of $r'!$ signed products
of $r'$ entries, at most one from the row $g_A^{\mathsf T}$. Multiplied
by $d^{r'}$, each product is an integer polynomial whose coefficient
absolute values sum to at most $C_0^{r'}$. Hence $\Delta=d^m$ and
$B=d^mm!\,C_0^m$ work.
\end{proof}

\begin{proof}[Proof of Proposition~\ref{prop:recourse-examples}]
(a) Here $\partial_z^2f=2v\ge0$. For $v>0$ the unique residual
minimizer is $z=0$, and at $z=1$ the gap is $v$ while the distance is
one; a bound $1\le\Gamma v^\theta$ for all $v\in(0,1]$ is impossible.
The Hessian of $f+\frac\alpha2v^2$ is
$\bigl[\begin{smallmatrix}\alpha&2z\\2z&2v\end{smallmatrix}\bigr]$, whose
determinant at $(0,\frac12)$ is $-1$. A function convex on the square
has positive semidefinite Hessian on the closed square, so no $\alpha$
works.

(b) The fibers are affine. For $v>0$ the coefficient of $z$ is
$-v^2<0$, so the unique residual minimizer is $z=1$ and the projected
value is $\gamma v-v^2=v(\gamma-v)>0$ on $(0,1]$, while $f_\gamma(0,z)=0$
for all $z$. Thus the optimal set is $\{0\}\times[0,1]$, $a=0$, $p=0$,
and $(b_j,1)\to(0,1)$. The Hessian of $f_\gamma+\frac\alpha2v^2$ is
$\bigl[\begin{smallmatrix}\alpha-2z&-2v\\-2v&0\end{smallmatrix}\bigr]$,
with determinant $-4v^2<0$ for $v>0$. For $f=-v^2z$, $\Lambda=0$ is a
valid coordinate bound, and when $\sigma>1$ the draws $\gamma\in(1,\sigma]$
have probability at least $(\sigma-1)/(2\sigma)-2/M>0$ under the grid
law for large $M$.
\end{proof}

\subsection{Proof of Lemma~\ref{lem:recourse-face}}
\label{app:recourse-face}

Fix $i$ and $\gamma_{-i}$, and let
\[
 V_i(s)=\min\Bigl\{f(v,z)+\sum_{l\ne i}\gamma_lv_l:\ (v,z)\in X,\ v_i=s\Bigr\},
 \qquad s\in[0,1].
\]
The $i$-th coordinates of the optimal cores of $f_{\gamma+\theta e_i}$
are exactly the minimizers of $V_i(s)+(\gamma_i+\theta)s$ on $[0,1]$. The
function $V_i$ is continuous, and Lipschitz with constant
$G_i=\sup_X\abs{\partial f/\partial v_i}$: moving the $i$-th coordinate
of a minimizer at $s'$ to $s$ stays in the product box. For every fixed
$(v_{-i},z)$, the function $s\mapsto f(v,z)-\bar\Lambda s^2/2$ is concave
on $[0,1]$, because $\partial_{v_i}^2f\le\Lambda\le\bar\Lambda$; as a
minimum of concave functions, $V_i-\bar\Lambda s^2/2$ is concave.

\emph{Monotonicity.} If $s_0$ minimizes $V_i+\gamma_is$ and $s_1$
minimizes $V_i+(\gamma_i-t)s$, adding the two optimality inequalities
gives $-ts_1\le-ts_0$, so $s_0\le s_1$.

\emph{Interior contacts.} If $s_0\in(0,1)$ minimizes $V_i+\vartheta s$,
then $V_i$ is differentiable at $s_0$ with $V_i'(s_0)=-\vartheta$.
Indeed, a supergradient $\zeta$ of the concave function
$V_i-\bar\Lambda s^2/2$ at $s_0$ gives
$V_i(s)\le V_i(s_0)+(\zeta+\bar\Lambda s_0)(s-s_0)+\frac{\bar\Lambda}2(s-s_0)^2$,
while minimality gives $V_i(s)\ge V_i(s_0)-\vartheta(s-s_0)$. Comparing
both sides of $s_0$ gives $\zeta+\bar\Lambda s_0=-\vartheta$, and $V_i$ is
squeezed between two functions with derivative $-\vartheta$ at $s_0$.
For two such points $s_0\le s_1$, concavity of $V_i-\bar\Lambda s^2/2$
gives $V_i'(s_1)-V_i'(s_0)\le\bar\Lambda(s_1-s_0)$.

(a) Let $b$ be an optimal core of $f_{\gamma-te_i}$ with
$b_i<t/\bar\Lambda\le\frac12$, and suppose an optimal core of $f_\gamma$
has $a_i>0$. By monotonicity $0<a_i\le b_i<\frac12$, so both are
interior contacts, with slopes $-\gamma_i$ and $-(\gamma_i-t)$. Then
$t=V_i'(b_i)-V_i'(a_i)\le\bar\Lambda(b_i-a_i)<\bar\Lambda b_i<t$, a
contradiction. The second rule follows by applying the first to
$s\mapsto V_i(1-s)$, which has the same semiconcavity.

(b) Let $\theta_i^0=\sup_{0<s\le1}(V_i(0)-V_i(s))/s$, finite because
$V_i$ is Lipschitz. The coordinate $0$ is optimal for the tilt
$\gamma_i$ exactly when $V_i(0)\le V_i(s)+\gamma_is$ for all $s$, that
is, when $\gamma_i\ge\theta_i^0$. If $\gamma_i-t>\theta_i^0$, then for
$s>0$, $V_i(s)+(\gamma_i-t)s>V_i(s)+\theta_i^0s\ge V_i(0)$, so every
optimal core of $f_{\gamma-te_i}$ has $i$-th coordinate $0$. Hence the
first situation in~(b) forces $\gamma_i\in[\theta_i^0,\theta_i^0+t]$.
With $\theta_i^1=\inf_{0\le s<1}(V_i(s)-V_i(1))/(1-s)$, the coordinate
$1$ is optimal exactly when $\gamma_i\le\theta_i^1$, and
$\gamma_i+t<\theta_i^1$ forces every optimal core of $f_{\gamma+te_i}$ to
have $i$-th coordinate $1$. Both thresholds depend only on
$\gamma_{-i}$. \qed

\begin{corollary}[Certified face rules]
\label{cor:recourse-face-rules}
Let $r_f=t/(8\bar\Lambda)$, and let $\hat b^{(i,\pm)}$ be rational
vectors within distance $r_f$ of the selected cores of
$f_{\gamma\pm te_i}$. Fix $v_i=0$ if $\hat b^{(i,-)}_i+r_f<t/\bar\Lambda$,
and fix $v_i=1$ if $\hat b^{(i,+)}_i-r_f>1-t/\bar\Lambda$. Then every fixed
value is correct for every optimal core of $f_\gamma$, no coordinate is
fixed twice, and under the grid law the probability that the selected
core $a(\gamma)$ has a coordinate in $\{0,1\}$ that is not fixed is at
most $kt/\sigma+2k/M$.
\end{corollary}

\begin{proof}
The rules imply the hypotheses of Lemma~\ref{lem:recourse-face}(a) for
the exact selected cores, so they are sound; two sound rules for the same
coordinate would contradict the existence of an optimal core. Suppose
$a_i=0$ and the first rule does not fire. Then the selected core $b$ of
$f_{\gamma-te_i}$ has $b_i\ge t/\bar\Lambda-2r_f>0$, so
$\gamma_i\in[\theta_i^0,\theta_i^0+t]$ by Lemma~\ref{lem:recourse-face}(b).
Conditionally on $\gamma_{-i}$, an interval of length $t$ contains at
most $t(M-1)/(2\sigma)+1$ grid points, so this has probability at most
$t/(2\sigma)+1/M$. The case $a_i=1$ is symmetric. A union bound over
$i$ and both endpoints gives the claim.
\end{proof}

\subsection{Proof of Lemma~\ref{lem:recourse-lattice}}
\label{app:recourse-lattice}

Every nonconstant monomial of degree at most two is $2$-Lipschitz on
$[0,1]^r$ in the maximum norm. Since both $a$ and $\hat a$ lie in
$[0,1]^r$, this gives
$\abs{T\chi_j(\hat a)-T\chi_j(a)}\le2T\norm{\hat a-a}_\infty\le\frac12$ and
$\abs{\omega_j-T\chi_j(a)}\le1$ for every $j$. The lattice $\mathcal L$
consists of the vectors $(h,h^{\mathsf T}\omega)$ with $h\in\Z^s$.

(a) If the test accepts, then
$\lambda_1(\mathcal L)\ge2^{-(s-1)/2}\norm{\mathbf b_1}>4sB$ by
Theorem~\ref{thm:recourse-lll}. Suppose $0<\norm h_\infty\le B$ and
$\abs{h^{\mathsf T}\chi(a)}\le sB/T$. Then
$\abs{h^{\mathsf T}\omega}\le\norm h_1+T\abs{h^{\mathsf T}\chi(a)}\le2sB$
and $\norm h\le sB$, so the nonzero lattice vector $(h,h^{\mathsf T}\omega)$
has length at most $\sqrt5sB<4sB$, a contradiction.

(b) Let $(h,h^{\mathsf T}\omega)\ne0$. If $\norm h>R_s$, its length
exceeds $R_s$. Otherwise $0<\norm h_\infty\le R_s$, and
$\abs{h^{\mathsf T}\omega}\ge T\abs{h^{\mathsf T}\chi(a)}-\norm h_1\ge
T\eta_0-sR_s>R_s$. Hence $\norm{\mathbf b_1}\ge\lambda_1(\mathcal L)>R_s=2^s4sB$,
and $\norm{\mathbf b_1}^2>2^{2s}(4sB)^2\ge2^{s-1}(4sB)^2$.

Constructing and rounding $\omega$ costs time polynomial in
$s+\log T+\langle\hat a\rangle$. The resulting basis entries have
$O(\log T)$ bits, so Theorem~\ref{thm:recourse-lll} gives the remaining
cost bound; the acceptance test is an exact integer comparison. \qed

Part~(b) must be stated with the larger height $R_s$: a short lattice
vector can have coefficients much larger than $B$, so margins over the
height-$B$ family alone would not imply acceptance.

\subsection{Proof of Proposition~\ref{prop:recourse-margin-probability}}
\label{app:recourse-margin}

\begin{lemma}[Quadratic sublevel sets]
\label{lem:recourse-sublevel}
For every nonzero integer polynomial $\psi$ of degree at most two in $r$
variables and every $0<\eta<1$,
$\operatorname{Leb}\{x\in[0,1]^r:\abs{\psi(x)}\le\eta\}\le8\sqrt\eta$.
\end{lemma}

\begin{proof}
For a univariate quadratic with leading coefficient $\ell\ne0$,
completing the square shows that $\{\abs{\cdot}\le\eta\}$ has total
length at most $2\sqrt{2\eta/\abs\ell}$; the two-interval case is
largest when the inner endpoints meet. If some square coefficient
$\psi_{ii}$ is nonzero, then $\abs{\psi_{ii}}\ge1$, and slicing along
$x_i$ with Fubini's theorem gives the bound $2\sqrt{2\eta}$. If all
square coefficients vanish but some $\psi_{ij}\ne0$, use the coordinates
$\varsigma=(x_i+x_j)/\sqrt2$ and $u=(x_i-x_j)/\sqrt2$. Then
$x_ix_j=(\varsigma^2-u^2)/2$, other terms are at most linear in
$\varsigma$, and every slice in $\varsigma$ has leading coefficient of
absolute value at least $\frac12$, hence length at most $4\sqrt\eta$;
the projection of the cube along $\varsigma$ has measure $\sqrt2$, which
gives $4\sqrt2\sqrt\eta$. If $\psi$ is a nonconstant linear polynomial,
slicing along a variable with a nonzero integer coefficient gives
$2\eta\le2\sqrt\eta$. A nonzero integer constant has empty sublevel set.
\end{proof}

\begin{lemma}[Contact slopes]
\label{lem:recourse-contact}
Let $\Phi:[0,1]^r\to\R$ be continuous with $\Phi-\frac U2\norm\cdot^2$
concave, and let $\mathcal C$ be the set of interior points $a$ at which
$\Phi(x)\ge\Phi(a)+s^{\mathsf T}(x-a)$ for all $x\in[0,1]^r$ and some $s$.
Then $s=s(a)$ is unique, $s$ is locally $U$-Lipschitz on $\mathcal C$,
and $\operatorname{Leb}^*(s(E\cap\mathcal C))\le U^r\operatorname{Leb}(E)$ for
every measurable $E$ in the interior of the cube.
\end{lemma}

\begin{proof}
As in the one-dimensional argument of Section~\ref{app:recourse-face},
a supergradient of the concave function at $a$ and the lower support
force $s(a)=\zeta+Ua$, and then
$0\le\Phi(x)-\Phi(a)-s(a)^{\mathsf T}(x-a)\le\frac U2\norm{x-a}^2$. Let
$a,a'\in\mathcal C$, $\varrho=\norm{a-a'}$, $\Delta s=s(a')-s(a)$, and
suppose the closed ball of radius $\varrho$ about $a$ lies in the cube.
For $\norm u=\varrho$, the lower support at $a'$ and the upper support at
$a$, evaluated at $a+u$, give
$\Delta s^{\mathsf T}u\le\Phi(a)-\Phi(a')-s(a')^{\mathsf T}(a-a')+\frac U2\varrho^2
\le U\varrho^2$, using the upper support at $a'$ evaluated at $a$.
Taking $u=\varrho\Delta s/\norm{\Delta s}$ gives
$\norm{\Delta s}\le U\varrho$. If $B(a_0,3\varrho_0)$ lies in the
interior, every pair of points of $\mathcal C\cap B(a_0,\varrho_0)$
satisfies the ball condition, so $s$ is $U$-Lipschitz there. Cover the
interior by countably many such balls $B_l$, disjointify
$E_l=E\cap B_l\setminus\bigcup_{l'<l}B_{l'}$, apply (F2) on each piece,
and sum: $\operatorname{Leb}^*(s(E\cap\mathcal C))\le
U^r\sum_l\operatorname{Leb}(E_l)=U^r\operatorname{Leb}(E)$.
\end{proof}

\begin{proof}[Proof of Proposition~\ref{prop:recourse-margin-probability}]
Identify the relative interior of $A$ with the interior of
$[0,1]^r$ through the free coordinates $x=v_F$, and put
$\Phi(x)=\min_{z\in[0,1]^m}f(v(x),z)$, where $v(x)$ has the fixed
coordinates of $A$. For every $z$, $x\mapsto f(v(x),z)-\frac U2\norm x^2$
is concave, so $\Phi-\frac U2\norm\cdot^2$ is concave. If a global
minimizer of $f_\gamma$ has core $a$ in the relative interior of $A$,
then $x_a=a_F$ minimizes $\Phi(x)+\gamma_F^{\mathsf T}x$ over the cube,
so $x_a\in\mathcal C$ with $s(x_a)=-\gamma_F$. With
$E=\{x:\abs{\psi(x)}\le\eta\}$ in the interior, the event therefore
implies $-\gamma_F\in s(E\cap\mathcal C)$. Conditionally on the noise of
the fixed coordinates, Lemmas~\ref{lem:recourse-contact}
and~\ref{lem:recourse-sublevel} bound its probability by
$U^r8\sqrt\eta/(2\sigma)^r$.

For the grid law, the event is defined by
\[
 \exists(a,z)\ \forall x':\ a\in\operatorname{relint}A\wedge z\in[0,1]^m
 \wedge\abs{\psi(a_F)}\le\eta\wedge
 \bigl[x'\notin X\vee f_\gamma(x')\ge f_\gamma(a,z)\bigr],
\]
a formula \eqref{eq:recourse-two-block} with two blocks of $n$
variables, $O(n)$ atoms, and degree at most three, in which $\psi$, $\eta$,
and all noise coordinates but one are real coefficients.
Lemmas~\ref{lem:recourse-sections} and~\ref{lem:recourse-replacement}
give the extra term $C_m/M$ with $C_m=2k(2H^2+1)$, independent of $\psi$
and $\eta$.
\end{proof}

\begin{corollary}[A uniform margin event]
\label{cor:recourse-margin-union}
Let $R\ge1$ be an integer, $s_*=\binom{k+2}2$, and
\[
 A_*=8\cdot3^k\max\{1,U/(2\sigma)\}^k(2R+1)^{s_*},\qquad
 C_{\rm tail}=3^k(2R+1)^{s_*}C_m .
\]
Call a draw \emph{bad} if, for some face $A$ with $r\ge1$ free
coordinates, some global minimizer of $f_\gamma$ has core $a$ in the
relative interior of $A$ and some nonzero integer polynomial $\psi$ of
degree at most two in the free coordinates, with coefficients of
absolute value at most $R$, has $\abs{\psi(a_F)}\le\eta$. Then
$\Pr_{\rm grid}\{\text{bad}\}\le A_*\sqrt\eta+C_{\rm tail}/M$.
\end{corollary}

\begin{proof}
There are at most $3^k$ faces and at most $(2R+1)^{s_*}$ polynomials per
face. Apply Proposition~\ref{prop:recourse-margin-probability} to each
pair and add. Faces with $r=0$ contribute nothing, since a nonzero
integer constant has absolute value at least $1>\eta$. The family is
used only in this union bound; no algorithm enumerates it.
\end{proof}

\subsection{Proofs of Proposition~\ref{prop:recourse-completion} and Theorem~\ref{thm:recourse-cubic}}
\label{app:recourse-cubic-proof}

\paragraph{Base choices.}
Let $\Xi\ge\max\{2,\Xi_0\}$ be as in \eqref{eq:recourse-core-params};
the same fallback factor $\Xi_0$ serves all instances below, because
they have the same format. Choose, in this order,
\begin{gather*}
 \bar\Lambda=\max\{1,\Lambda\},\qquad
 t=\min\Bigl\{\frac12,\frac{\sigma}{16k\Xi}\Bigr\},\qquad
 r_f=\frac t{8\bar\Lambda},\\
 \Delta,B\ \text{from Lemma~\ref{lem:recourse-minor-heights}},\qquad
 s_*=\binom{k+2}2,\qquad R=2^{s_*}4s_*B,\\
 U\ge\max\Bigl\{1,\sum_{i,j}\sup_X\abs{\partial_{v_i}\partial_{v_j}f}\Bigr\}
 \ \text{by coefficient sums},\qquad
 \eta=\min\Bigl\{\frac14,\frac1{(16A_*\Xi)^2}\Bigr\},\\
 T=\text{the least power of two with }T>(s_*+1)R/\eta,
\end{gather*}
with $A_*$ from Corollary~\ref{cor:recourse-margin-union}. Finally let
$M$ be the least power of two that satisfies \eqref{eq:recourse-core-params}
for each of the $2k+1$ base objectives $f$ and $f\pm tv_i$ and also
$(2k+C_{\rm tail})/M\le1/(8\Xi)$. The base objectives $f\pm tv_i$ have
the same residual convexity and coordinate bound $\Lambda$ as $f$, and
input length polynomial in $L$. Every requirement on $M$ is a lower
bound, so one grid serves all of them. All quantities have bit length
polynomial in $L$ and precede sampling. On a draw $\gamma$, the core
evaluator of Theorem~\ref{thm:recourse-core} for the base objective
$f\pm tv_i$ approximates the selected core of $f_{\gamma\pm te_i}$; the
draw is not resampled.

\paragraph{Base stage.}
(1) For each $i$, query the core evaluators of $f_{\gamma\pm te_i}$ with
accuracy $r_f$ and apply the rules of
Corollary~\ref{cor:recourse-face-rules}. Let $A$ be the face defined by
the fixed coordinates, $F$ its free coordinates, and $r=\abs F$.
(2) If $r=0$, accept with $\delta=\frac12$ and $\mu=\min\{1,1/\Delta\}$.
Otherwise query the core evaluator of $f_\gamma$ with accuracy $1/(4T)$,
clip the free coordinates of the answer to $[0,1]$ to obtain $\hat a$,
and run the test of Lemma~\ref{lem:recourse-lattice} with these $B$ and
$T$ and $s=\binom{r+2}2$. If it rejects, reject. If it accepts, put
$\xi=sB/T$, $\delta=\min\{\frac12,\xi\}$, and $\mu=\min\{1,\xi/\Delta\}$.
(3) On acceptance compute $\Gamma$ from
Proposition~\ref{prop:recourse-fiber}(c) for $A$, $\delta$, and $\mu$.
The outcome is a deterministic function of the draw.

\paragraph{Acceptance is sound.}
By Corollary~\ref{cor:recourse-face-rules}, $a(\gamma)\in A$. If $r=0$,
then $a(\gamma)$ is the vertex $A$; the minors of $J_A$ are rational
constants with denominators dividing $\Delta$, so the nonzero ones have
absolute value at least $1/\Delta\ge\mu$, and the margin condition on
free coordinates is vacuous. If $r\ge1$ and the test accepts, then
$\hat a\in[0,1]^r$ by clipping, and $\norm{\hat a-a_F}_\infty\le1/(4T)$
because clipping to a box containing $a_F$ does not increase any
coordinate error. So Lemma~\ref{lem:recourse-lattice}(a) applies: every
nonzero integer quadratic of height at most $B$ in the free coordinates
has absolute value greater than $\xi$ at $a_F$. These
include $v_i$ and $1-v_i$ for $i\in F$, so the free coordinates of $a$
lie in $(\xi,1-\xi)$, and $a$ lies in the relative interior of $A$ with
margin $\delta$. By Lemma~\ref{lem:recourse-minor-heights}, each square
minor of $J_A$ is either identically zero on $A$, hence zero at $a$, or
equal to $\Delta^{-1}$ times a nonzero integer quadratic of height at
most $B$, hence of absolute value greater than $\xi/\Delta\ge\mu$ at $a$.
Proposition~\ref{prop:recourse-fiber}(c) therefore gives
\begin{equation}
 \dist(z,S(\gamma))\le\Gamma\bigl(f(a,z)-\min f(a,\cdot)\bigr)^{1/4}
 \quad\text{for } z\in[0,1]^m \text{ with gap at most } 1,
 \label{eq:recourse-accepted-bound}
\end{equation}
since $S(\gamma)=S_a$. No probability was used. The supplied margins
$\delta$ and $\mu$ are rationals with
$\langle\delta\rangle+\langle\mu\rangle=O(\log s+\log B+\log T+\log\Delta)
=\poly(L)$, so $\Gamma$ is computed in time $\poly(L)$.

\paragraph{Completion after acceptance.}
For a query $q$ put $\epsilon=2^{-q}$, $e=\epsilon/2$,
$R_z=\max\{1,\lceil\sqrt m\rceil\}$, a rational
$G\ge\max\{1,\sup_X\norm{\nabla_vf}\}$, and
\[
 \tau=\frac{e^6}{1024R_z^4\Gamma^4},\qquad
 \eta_q=\frac{\tau e^2}8,\qquad
 d_q=\min\Bigl\{e,\frac{\tau e^2}{16G}\Bigr\}.
\]
Query the core evaluator of $f_\gamma$ with accuracy $d_q$, set the
coordinates fixed by $A$ to their exact values, and clip the others to
$[0,1]$. This gives a rational $b\in A$ with $\norm{b-a}\le d_q$, since
$a\in A\cap[0,1]^k$. Apply \eqref{eq:recourse-cv} to the rational convex
cubic $Q_b(z)=f(b,z)+\tau\norm z^2$ on $[0,1]^m$ with accuracy $\eta_q$,
and return $(b,\hat z)$. Since $\abs{f(b,z)-f(a,z)}\le G\norm{b-a}$ for
all $z$, the function $Q_a(z)=f(a,z)+\tau\norm z^2$ satisfies
$Q_a(\hat z)-\min Q_a\le\eta_q+2Gd_q\le\tau e^2/4$.
Lemma~\ref{lem:recourse-selection}, with $\varphi=f(a,\cdot)$ on
$[0,1]^m$, the bound \eqref{eq:recourse-accepted-bound}, and $D_e=4$,
gives $\norm{\hat z-p}\le e$. Hence
$\norm{(b,\hat z)-(a,p)}\le\sqrt2\,e<\epsilon$. The core query has
precision $\log(1/d_q)\le8q+\poly(L)$, and the convex solve has data of
bit length $\poly(L+q)$ and tolerance $\langle\eta_q\rangle\le8q+\poly(L)$.
This proves
Proposition~\ref{prop:recourse-completion}.

\paragraph{Rejection is rare.}
Suppose that no coordinate of $a(\gamma)$ in $\{0,1\}$ is missed by the
face rules and that the draw is not bad in the sense of
Corollary~\ref{cor:recourse-margin-union}. Since the rules are sound,
$A$ is then the smallest face containing $a$, and $a$ lies in its
relative interior. If $r=0$ the stage accepts. Otherwise every nonzero
integer quadratic of height at most $R$ in the free coordinates has
absolute value greater than $\eta$ at $a_F$. Because $s\le s_*$, we have
$R_s=2^s4sB\le R$ and $T\eta>(s_*+1)R\ge(s+1)R_s$, so
Lemma~\ref{lem:recourse-lattice}(b) shows that the test accepts. By
Corollaries~\ref{cor:recourse-face-rules}
and~\ref{cor:recourse-margin-union},
\[
 \Pr\{\text{reject}\}\le\frac{kt}\sigma+\frac{2k}M+A_*\sqrt\eta
 +\frac{C_{\rm tail}}M\le\frac1{16\Xi}+\frac1{16\Xi}+\frac1{8\Xi}
 =\frac1{4\Xi}.
\]

\paragraph{The same selector on rejection.}
On rejection, return the output of Lemma~\ref{lem:recourse-fallback} for
$f_\gamma$ with precision $q+\lceil\log_2n\rceil$; its $\infty$-accuracy
is at most $2^{-q}/n$, so its Euclidean accuracy is at most $2^{-q}$.
Its singleton formulas
\eqref{eq:recourse-singleton} define $x^\sharp(\gamma)$ itself: global
optimality, the lexicographically least core, and the least residual
norm within that core. So the ordinary and fallback branches approximate
the same point at every precision, including on draws with several
optimal cores. The cost is at most $\Xi P(L+q)$.

\paragraph{Work.}
Let $\Omega_0$ and $\Omega_{i,\pm}$ be the random factors of
Theorem~\ref{thm:recourse-core} for the base objectives $f$ and
$f\pm tv_i$. The base stage makes $2k+1$ core queries of precision
$\poly(L)$, one LLL reduction, and polynomial arithmetic for $\Gamma$; a
completion makes one core query of precision $\poly(L)+O(q)$ and one
convex solve. Every query therefore costs at most $\Omega(\gamma)P(L+q)$
for a fixed polynomial $P$, where
\[
 \Omega=\Omega_0+\sum_{i,\pm}\Omega_{i,\pm}+\Xi\,\mathbf1\{\text{reject}\}.
\]
By Theorem~\ref{thm:recourse-core} and the rejection bound,
$\E\,\Omega\le(2k+1)k^{O(k)}(1+\Lambda/\sigma)^k\poly(L)+\frac14$. No
independence between the factors and the rejection event is used. The
numbers $U$, $B$, $T$, and $\Gamma$ enter only through their bit lengths;
the numerical factor remains $(1+\Lambda/\sigma)^k$. This proves
Theorem~\ref{thm:recourse-cubic}. \qed

\begin{proof}[Proof of Remark~\ref{rem:recourse-value-certificate}]
Let $x_q$ be the returned point at accuracy $2^{-q-1}/G$ and $\ell_{q+1}$
the lower bound from Theorem~\ref{thm:recourse-core}. Then
$\ell_{q+1}\le f_\gamma^*\le f_\gamma(x_q)\le f_\gamma^*+2^{-q-1}\le
\ell_{q+1}+2^{-q}$, because $X$ is convex and $\norm{\nabla f_\gamma}\le G$
on $X$.
\end{proof}

\begin{proof}[Proof of Remark~\ref{rem:recourse-quartic}]
For $f(v,y)=v^2/2+F(y)$ and every draw $\gamma$, the objective separates:
the core part $v^2/2+\gamma v$ has a unique minimizer on $[0,1]$, and the
residual part has the unique minimizer $y^*$ of $F$. So
$x^\sharp(\gamma)=(a(\gamma),y^*)$ on every draw. The residual fibers are
convex because $F$ is convex on its box, and $\partial_v^2f=1$. A query
with $q=2$ returns a point within $1/4$ of $x^\sharp(\gamma)$, which
locates the designated coordinate of $y^*$ to within $1/4$ and decides
the source instance, as in Theorem~\ref{thm:points-quartic-lower}. An
evaluator that is correct on every draw with polynomial expected work
would thus give a Las Vegas algorithm with expected polynomial running
time: since $M$ is a power of two, a draw uses $\log_2M$ fair coins, and
with $\sigma=1$ the factor $(1+\Lambda/\sigma)^k$ equals $2$.
\end{proof}
